\chapter{PROOFS AND SUPPLEMENTARY DETAILS FOR CHAPTER~\ref{ch:pomdp-control}}
\label{app:pomdp-control}

\aigen{This appendix contains the complete belief-tree planner pseudocode (Algorithms~\ref{alg:mctq} and \ref{alg:simulate}), the practical state-rollout deep RL instantiation (Algorithm~\ref{alg:pf_mcts_practical}), the proofs of the four-way regret decomposition and special-case QMDP structural bias bounds, environment specifications and posterior diagnostics for Double-Integrator LQR and HeartSteps, hyperparameter tables, and the camera-ready validation and ablation studies for Chapter~\ref{ch:pomdp-control}.}


\section[Finite-Depth Belief-Tree Planning Subroutine]{\aigen{Finite-Depth Belief-Tree Planning Subroutine (Algorithm~\ref{alg:mctq})}}
\label{app:mctq}

We describe the full finite-depth Monte Carlo belief-tree planning subroutine with QMDP leaves (\aigen{Algorithm~\ref{alg:mctq}}) referenced in Algorithm~\ref{alg:pf_mcts}.
\aigen{Algorithm~\ref{alg:mctq}} estimates the depth-$d$ belief-tree action-value $Q_{\theta,d}(b^J,a)$ defined in~\eqref{eq:Qd_def} by Monte Carlo tree search on the particle belief.
It is inspired by the POMCP framework~\citep{silver2010monte}, with three departures: internal actions are drawn uniformly at random rather than by UCB1, every simulation descends to the full depth $d$, and values are computed by a bottom-up dynamic programming (DP) backup after all simulations complete.

\paragraph{Overview.}
Given the current particle belief $b^J=\frac{1}{J}\sum_{j=1}^J\delta_{s_j}$, model $\theta$, search depth $d$, simulation budget $N$ per root action, and a precomputed fully observable Q-function $Q_\theta^{\star,S}$ (e.g., from value iteration), \aigen{Algorithm~\ref{alg:mctq}} estimates $Q_{\theta,d}(b^J,a)$ for each root action $a\in\gA$ separately (round-robin), so that every root action receives exactly $N$ simulations.
For each root action $a$, each simulation samples a particle $s_j$ from $b^J$, takes the fixed root action $a$ to obtain a child state $s'\sim P_\theta(\cdot\mid s_j,a)$ and observation $o\sim Z_\theta(\cdot\mid s')$, and then descends the search tree $T_a$ from the child node $(a,o)$ down to depth $d$, selecting actions uniformly at random at each internal node.
Nodes are created the first time a simulation passes through them, and every node on the path records the simulated state.
After all $N$ simulations complete, a bottom-up DP backup computes value estimates: at each depth-$d$ leaf, the QMDP value of its accumulated particles; at each internal node, $\hat V(v)=\max_{a'}\hat Q(v,a')$, where $\hat Q(v,a')$ is the empirical mean of $r+\gamma\hat V(\text{child})$ over the simulations that selected $a'$ at $v$.
The root estimate $\widehat Q(b^J,a)$ is the same empirical mean at the root.

Like POMCP~\citep{silver2010monte}, \aigen{Algorithm~\ref{alg:mctq}} indexes its search trees by action-observation histories $h=(a_1,o_1,\ldots,a_k,o_k)$ rather than by explicit beliefs, and represents the belief at each node implicitly by the collection of states simulated through that node.
The round-robin over root actions is needed because the regret analysis (Lemma~\ref{lem:delta_bound_main}) requires uniformly good estimates $\widehat Q(b^J,a)$ for \emph{all} $a\in\gA$.
The analysis (Proposition~\ref{prop:mc_concentration}) uses two properties of the internal action-selection rule $\pi$, i.e., the rule that picks actions at nodes below the root within one planning call:
\begin{enumerate}
\item[(S1)] \emph{Non-adaptivity.} $\pi$ may depend on the node's history, on $\theta$, on the root particle set $b^J$, and on fresh randomness, but not on any statistic accumulated by the current tree.
\item[(S2)] \emph{Full support.} $\pi(a'\mid\cdot)\ge\pi_{\min}>0$ for every $a'\in\gA$.
\end{enumerate}
Uniform selection satisfies both with $\pi_{\min}=1/|\gA|$.
So does a learned proposal that is frozen during the planning call and mixed with the uniform distribution.
UCB1 violates (S1), because its choice at a node reads the visit counts and value estimates accumulated by the current tree.
Uniform selection is therefore a proof device, not a performance claim.
Together with the absence of an early return (every simulation descends to depth $d$), (S1) makes the path of each simulation a function of its own random seed alone (Lemma~\ref{lem:path_decoupling}).
The simulation seeds are i.i.d.\ given $b^J$, but the simulation returns are not, because the leaf evaluator pools particles across simulations and the DP backup couples them.
The analysis therefore treats $\widehat Q(b^J,a)$ as a bounded-difference function of the i.i.d.\ seeds rather than as an average of independent returns.

\paragraph{Tree structure and belief particles.}
Each search tree $T_a$ (one per root action $a$) has the particle belief $b^J$ at its root (depth $0$), the root action fixed to $a$, internal nodes at depths $1,\ldots,d-1$, and leaves at depth $d$.
Each node $v$ is identified with the action-observation history that leads to it from the root, so a given history always refers to the same node.
A child of node $v$ under action $a'$ and observation $o$ is denoted $\textsc{Child}(v,a',o)$; it is created on demand when first visited.

Each node $v$ maintains a \emph{particle set} $B(v)$: the collection of all states that have been simulated through $v$ across all simulations.
The empirical distribution $\hat b_v:=\frac{1}{|B(v)|}\sum_{s\in B(v)}\delta_s$ is the particle approximation of the belief at $v$, and $N(v):=|B(v)|$ is the number of simulations that have visited $v$.
In addition, each internal node stores:
\begin{itemize}
\item $N(v,a')$: the number of simulations that selected action $a'$ at $v$;
\item $\bar r(v,a')$: the running mean of the immediate rewards $r(s,a')$ over those simulations.
\end{itemize}
These statistics are accumulated during the simulation phase and used in the post-simulation DP backup.

\paragraph{QMDP leaf evaluation.}
At depth-$d$ leaf nodes (remaining depth $\ell=0$), the QMDP value of the accumulated belief is
\begin{equation}\label{eq:qmdp_leaf_eval}
\bar V_\theta^{\star,S}(B(v)):=\max_{a'\in\gA}\frac{1}{|B(v)|}\sum_{s\in B(v)}Q_\theta^{\star,S}(s,a').
\end{equation}
This takes the max \emph{outside} the average over states, choosing a single action for all particles---exactly the QMDP approximation $\bar V_\theta^{\star,S}$ from~\eqref{eq:U0_def}. This is distinct from the per-state evaluation $\frac{1}{|B(v)|}\sum_s\max_{a'}Q_\theta^{\star,S}(s,a')$, which would compute the linear surrogate $\widetilde V_\theta^{\star}$ instead. Using~\eqref{eq:qmdp_leaf_eval} ensures the algorithm matches the theoretical target $\bar V_{\theta,d}^{\star,S}=\gT_\theta^d\bar V_\theta^{\star,S}$ analyzed in Section~\ref{sec:analysis}.

\paragraph{Simulation procedure.}
For each root action $a$, each of the $N$ simulations proceeds as follows.
A particle $s_j$ is sampled uniformly from $b^J$, and the root transition is simulated: $s'\sim P_\theta(\cdot\mid s_j,a)$, $o\sim Z_\theta(\cdot\mid s')$, with immediate reward $r(s_j,a)$.
Starting from the child node $\textsc{Child}(\mathrm{root},a,o)$ with remaining depth $\ell=d-1$, the \textsc{Simulate} subroutine (Algorithm~\ref{alg:simulate}) descends the tree recursively:
\begin{enumerate}
\item \emph{Record.} Add the current state $s$ to $B(v)$.
\item \emph{QMDP leaf ($\ell=0$).} Stop. The leaf is evaluated by~\eqref{eq:qmdp_leaf_eval} on its final particle set once all simulations are complete.
\item \emph{Expansion (new internal node).} If $v$ is not yet expanded, initialize $N(v,a')\gets0$ and $\bar r(v,a')\gets0$ for all $a'\in\gA$, and continue. There is no early return: a simulation that creates a node keeps descending.
\item \emph{Uniform selection.} Select
\begin{equation}\label{eq:uniform_action}
a'\sim\mathrm{Unif}(\gA).
\end{equation}
\item \emph{Simulate and recurse.} Draw $s''\sim P_\theta(\cdot\mid s,a')$ and $o'\sim Z_\theta(\cdot\mid s'')$, update $N(v,a')$ and $\bar r(v,a')$ with the reward $r(s,a')$, and recurse into $\textsc{Child}(v,a',o')$ with remaining depth $\ell-1$.
\end{enumerate}
Every simulation therefore visits exactly one node at each depth $1,\ldots,d$.
After all $N$ simulations for root action $a$, the DP backup (Phase~2 of Algorithm~\ref{alg:mctq}) computes, from the leaves upward,
\begin{equation}\label{eq:dp_backup}
\begin{aligned}
\hat V(v)&=\bar V_\theta^{\star,S}(B(v)) &&\text{(leaf $v$)},\\
\hat Q(v,a')&=\bar r(v,a')+\gamma\sum_{w\in\mathrm{Ch}(v,a')}\frac{N(w)}{N(v,a')}\,\hat V(w) &&\text{(internal $v$, $N(v,a')>0$)},\\
\hat V(v)&=\max_{a':\,N(v,a')>0}\hat Q(v,a') &&\text{(internal $v$)},
\end{aligned}
\end{equation}
where $\mathrm{Ch}(v,a')$ is the set of children of $v$ under action $a'$.
Equivalently, $\hat Q(v,a')$ is the average of $r(s,a')+\gamma\hat V(\text{child})$ over the simulations that selected $a'$ at $v$.
The root estimate is
\[
\widehat Q(b^J,a)=\frac{1}{N}\sum_{n=1}^N\bigl[r(s_{j_n},a)+\gamma\,\hat V(v_{n,1})\bigr],
\]
where $s_{j_n}$ is the particle sampled by simulation $n$ and $v_{n,1}$ is the depth-$1$ node it enters.
When $d=0$ no tree is built and $\widehat Q(b^J,a)=\frac1J\sum_{j=1}^J Q_\theta^{\star,S}(s_j,a)$.

The full procedure is given in Algorithms~\ref{alg:mctq} and~\ref{alg:simulate}.

\begin{algorithm}[t]
\caption[Finite-depth Monte Carlo belief-tree planning with QMDP leaf bootstrapping]{\aigen{Finite-Depth Monte Carlo Belief-Tree Planning with QMDP Leaf Bootstrapping}}
\label{alg:mctq}
\begin{algorithmic}[1]
\Require model $\theta$; particle belief $b^J=\frac{1}{J}\sum_{j=1}^J\delta_{s_j}$; depth $d$; simulations per action $N$; MDP Q-function $Q_\theta^{\star,S}$
\If{$d=0$} \Comment{No tree: exact QMDP}
    \State\Return $\bigl\{\frac1J\sum_{j=1}^J Q_\theta^{\star,S}(s_j,a)\bigr\}_{a\in\gA}$
\EndIf
\For{each action $a\in\gA$} \Comment{Round-robin: $N$ simulations per root action}
    \State Initialize empty search tree $T_a$
    \For{$n=1,\ldots,N$} \Comment{Phase 1: tree construction}
        \State Sample $j_n\sim\mathrm{Unif}([J])$
        \State Sample $s'\sim P_\theta(\cdot\mid s_{j_n},a)$, $o\sim Z_\theta(\cdot\mid s')$
        \State $v_{n,1}\gets\textsc{Child}(\mathrm{root}(T_a),a,o)$ \Comment{Child of root for $(a,o)$}
        \State $\textsc{Simulate}(s',\,d-1,\,\theta,\,v_{n,1})$ \Comment{Alg.~\ref{alg:simulate}}
    \EndFor
    \State \Comment{Phase 2: DP backup~\eqref{eq:dp_backup}}
    \State $\hat V(v)\gets\max_{a'\in\gA}\frac{1}{|B(v)|}\sum_{s\in B(v)}Q_\theta^{\star,S}(s,a')$ at each leaf $v$ of $T_a$ \Comment{Pooled QMDP leaf}
    \For{each internal node $v$ of $T_a$, from depth $d-1$ up to depth $1$}
        \State $\hat Q(v,a')\gets\bar r(v,a')+\gamma\sum_{w\in\mathrm{Ch}(v,a')}\frac{N(w)}{N(v,a')}\hat V(w)$ for each $a'$ with $N(v,a')>0$
        \State $\hat V(v)\gets\max_{a':\,N(v,a')>0}\hat Q(v,a')$
    \EndFor
    \State $\widehat Q(b^J,a)\gets\frac{1}{N}\sum_{n=1}^N\bigl[r(s_{j_n},a)+\gamma\,\hat V(v_{n,1})\bigr]$
\EndFor
\State\Return $\{\widehat Q(b^J,a)\}_{a\in\gA}$
\end{algorithmic}
\end{algorithm}

\begin{algorithm}[t]
\caption[\textsc{Simulate}: recursive uniform-exploration descent to full depth]{\textsc{Simulate}: Recursive uniform-exploration descent to full depth}
\label{alg:simulate}
\begin{algorithmic}[1]
\Require state $s$; remaining depth $\ell$; model $\theta$; tree node $v$
\State $B(v)\gets B(v)\cup\{s\}$ \Comment{Record state; $N(v)=|B(v)|$}
\If{$\ell=0$} \Comment{Leaf: evaluated in Phase 2 of Alg.~\ref{alg:mctq}}
    \State\Return
\EndIf
\If{$v$ is not yet expanded} \Comment{Expansion: initialize and continue (no early return)}
    \State Expand $v$: $N(v,a')\gets 0$, $\bar r(v,a')\gets 0$ for all $a'\in\gA$
\EndIf
\State Sample $a'\sim\mathrm{Unif}(\gA)$ \Comment{Uniform, reads no tree statistic}
\State Sample $s'\sim P_\theta(\cdot\mid s,a')$, $o\sim Z_\theta(\cdot\mid s')$
\State $N(v,a')\gets N(v,a')+1$; $\bar r(v,a')\gets\bar r(v,a')+\frac{r(s,a')-\bar r(v,a')}{N(v,a')}$ \Comment{Incremental mean}
\State $\textsc{Simulate}(s',\,\ell-1,\,\theta,\,\textsc{Child}(v,a',o))$
\end{algorithmic}
\end{algorithm}

\paragraph{Convergence.}
As $N\to\infty$, two things happen.
First, by Lemma~\ref{lem:path_decoupling} the events ``simulation $n$ reaches $v$ and selects $a'$'' are i.i.d.\ across $n$ with positive probability for every reachable node $v$ and every $a'$, so every such node and action is visited infinitely often.
Second, the states entering a node $v$ are i.i.d.\ draws from the (binned) conditional belief $b_v$ at $v$, so $\hat b_v\to b_v$ by the law of large numbers.
At each depth-$d$ leaf, $\bar V_\theta^{\star,S}(B(v))\to\bar V_\theta^{\star,S}(b_v)$, and the DP backup then gives, level by level, $\hat Q(v,a')\to Q_{\theta,\ell}^{\phi}(b_v,a')$ and $\hat V(v)\to \max_{a'}Q_{\theta,\ell}^{\phi}(b_v,a')=\bar V_{\theta,\ell}^{\star,S}(b_v)$, where $\ell$ is the remaining depth at $v$.
Since the root action is fixed and states are sampled from $b^J$, the backed-up estimate $\widehat Q(b^J,a)\to Q_{\theta,d}^{\phi}(b^J,a)$ for each $a\in\gA$.
Proposition~\ref{prop:mc_concentration} quantifies the rate.

When $d=0$, no tree search is performed and the planner returns $\max_a\frac{1}{J}\sum_j Q_\theta^{\star,S}(s_j,a)$---exactly the QMDP value $\bar V_\theta^{\star,S}(b^J)$.
When $d\to\infty$, the tree-search value converges to $Q_\theta^{\star,B}(b^J,a)$.

\paragraph{Computational cost.}
For each root action, each of the $N$ simulations traverses a path of exactly $d$ nodes in the tree, requiring $O(d)$ generative model calls and $O(1)$ work for uniform action sampling at each internal node.
The post-simulation DP backup adds $O(|\gA|)$ work per node but is performed once, not per simulation.
The QMDP leaf evaluation~\eqref{eq:qmdp_leaf_eval} can be maintained incrementally by storing running sums $\sum_{s\in B(v)}Q_\theta^{\star,S}(s,a')$ for each action $a'$, reducing the per-simulation leaf cost to $O(|\gA|)$.
The total cost per planning step is $O(|\gA|^2 Nd)$ model evaluations across all root actions.

\paragraph{Relation to the error analysis.}
The Monte Carlo error $\epsilon_{\mathrm{mc}}(t,d)=\sup_a|\widehat Q_t(b_t^J,a)-Q_{\theta,d}(b_t^J,a)|$ arises from the finite simulation budget $N$ per root action.
The simulation seeds are i.i.d.\ given $b_t^J$, and under (S1) with no early return each simulation's path is a function of its own seed (Lemma~\ref{lem:path_decoupling}).
The estimate $\widehat Q_t(b_t^J,a)$ is then a bounded-difference function of the seeds, so the Efron--Stein and McDiarmid inequalities control its fluctuation at rate $O(\Lambda_d/\sqrt N)$, while the DP-backup bias vanishes at rate $O(\Phi_d/\sqrt{N})$ (Proposition~\ref{prop:mc_concentration}).
The simulation returns themselves are not independent, and the argument does not require them to be.
The proof is written for uniform selection; for a rule satisfying (S1)--(S2), it goes through with $|\gA|$ replaced by $1/\pi_{\min}$ wherever it bounds the fraction of a node's visits that select a given action.



\section{Practical Instantiation and Benchmark Specifications}
\label{ch08:app:experiments}
\label{alg:pf_mcts_practical}
\label{app:lqr}
\label{fig:lqr-snapshot}
\label{fig:lqr-posterior}
\label{fig:lqr-belief-mean}
\label{app:heartsteps}
\label{fig:heartsteps-all}
\label{fig:heartsteps-snapshot}
\label{fig:heartsteps-posterior}

\aigen{Extended related work on belief-space planning and particle filtering, the practical deep reinforcement learning instantiation (Algorithm~\ref{alg:pf_mcts_practical}), and full environment specifications, posterior diagnostics, and hyperparameter tables for the Double-Integrator LQR, HeartSteps, and MuJoCo benchmarks appear in Appendices~A and~C of \citet{tan2026pomdp}.}

\section{Detailed Error Analysis and Proofs}
\label{app:proofs}

Throughout, $|r(s,a)|\le R_{\max}$ and $V_{\max}:=R_{\max}/(1-\gamma)$. We define the per-round error quantities:
\begin{align*}
\epsilon_{\mathrm{mc}}(t,d)&:=\sup_{a\in\gA}\bigl|\widehat Q_t(b_t^J,a)-Q_{\theta,d}(b_t^J,a)\bigr|, &\text{(Monte Carlo error)}\\
\epsilon_{\mathrm{struct}}(t,d)&:=\sup_{a\in\gA}\bigl|Q_{\theta}^{\star,B}(b_t^J,a)-Q_{\theta,d}(b_t^J,a)\bigr|. &\text{(structural error)}
\end{align*}


%% ===================================================================
\subsection{Structural Error of Leaf Value Bootstrapping}
\label{app:structural_proofs}

Define the \emph{linear QMDP surrogate}
\begin{equation}
\widetilde V_\theta^{\star}(b):=\int_{\gS}V_\theta^{\star,S}(s)\,b(\mathrm ds).
\label{eq:linear_surrogate}
\end{equation}
By Jensen's inequality, $\widetilde V_\theta^{\star}(b)\ge \bar V_\theta^{\star,S}(b)$ and the gap $\Xi_\theta := \sup_b [\widetilde V_\theta^{\star}(b)-\bar V_\theta^{\star,S}(b)] \ge 0$.

\begin{prop}[Structural Error of the QMDP Leaf Value Approximation]
\label{prop:structural_error}
For every $\theta\in\Theta$ and $d\ge 0$,
\begin{equation}
\|\bar V_{\theta, d}^{\star,S}-V_{\theta}^{\star,B}\|_\infty
\;\le\;
\frac{\gamma^{d+1}}{1-\gamma}\,\Xi_\theta.
\label{eq:qmdp_gap}
\end{equation}
\end{prop}

In particular, $\epsilon_{\mathrm{struct}}(u,d)\le\frac{\gamma^{d+1}}{1-\gamma}\Xi_{\bar\theta_u}$.

\begin{proof}[Proof of Proposition~\ref{prop:structural_error}]
Write
\begin{align*}
Q_{\theta}^{\star,B}(b,a)
&=r(b,a)+\gamma\!\int_{\gO}V_{\theta}^{\star,B}(\tau_\theta(b,a,o))\,p_\theta(o\mid b,a)\,\mathrm do,\\
Q_{\theta,0}(b,a)
&=r(b,a)+\gamma\!\int_{\gS}V_\theta^{\star,S}(s')\,\tilde b_\theta(\mathrm ds'\mid b,a)
=r(b,a)+\gamma\,\widetilde V_\theta^{\star}(\tilde b_\theta(\cdot\mid b,a)).
\end{align*}
By Bayes consistency and linearity of $\widetilde V_\theta^{\star}$:
\[
\widetilde V_\theta^{\star}(\tilde b_\theta(\cdot\mid b,a))
=\int_{\gO}\widetilde V_\theta^{\star}(\tau_\theta(b,a,o))\,p_\theta(o\mid b,a)\,\mathrm do.
\]
Subtracting and using $\widetilde V_\theta^{\star}=\bar V_\theta^{\star,S}+(\widetilde V_\theta^{\star}-\bar V_\theta^{\star,S})$ with $0\le \widetilde V_\theta^{\star}-\bar V_\theta^{\star,S}\le\Xi_\theta$:
\[
\|Q_\theta^{\star,B}-Q_{\theta,0}\|_\infty
\;\le\;
\gamma\|V_\theta^{\star,B}-\bar V_\theta^{\star,S}\|_\infty+\gamma\,\Xi_\theta.
\]
Since $\|V_\theta^{\star,B}-\bar V_\theta^{\star,S}\|_\infty=\sup_b|\max_a Q_\theta^{\star,B}(b,a)-\max_a Q_{\theta,0}(b,a)|\le\|Q_\theta^{\star,B}-Q_{\theta,0}\|_\infty$ (using $|\max f-\max g|\le\max|f-g|$), rearranging gives $\|V_\theta^{\star,B}-\bar V_\theta^{\star,S}\|_\infty\le\frac{\gamma}{1-\gamma}\Xi_\theta$ (the $d=0$ case). For general $d\ge 0$: $\bar V_{\theta,d}^{\star,S}=\gT_\theta^d\bar V_\theta^{\star,S}$ and $V_\theta^{\star,B}=\gT_\theta V_\theta^{\star,B}$ (fixed point), so
\[
\|\bar V_{\theta,d}^{\star,S}-V_\theta^{\star,B}\|_\infty
=\|\gT_\theta^d\bar V_\theta^{\star,S}-\gT_\theta^d V_\theta^{\star,B}\|_\infty
\le\gamma^d\|\bar V_\theta^{\star,S}-V_\theta^{\star,B}\|_\infty
\le\frac{\gamma^{d+1}}{1-\gamma}\Xi_\theta.\qedhere
\]
\end{proof}


%% ===================================================================
\subsection{Planning and State-Identification Error}
\label{app:planning_state_proofs}

\begin{defn}[Q-belief distance]\label{def:dQ}
For $b,b'\in\Delta(\gS)$, define
$d_Q(b,b'):=\sup_{\theta\in\Theta,\,a\in\gA}|Q_\theta^{\star,B}(b,a)-Q_\theta^{\star,B}(b',a)|$.
\end{defn}

\begin{lem}[$Q^{\star,B}$-difference via test functions]\label{lem:Q_lipschitz}
Fix $\theta\in\Theta$, $a\in\gA$, and truncation depth $H\ge 1$. Recall the Bayesian update $\tau(b,a,o)\in\Delta(\gS)$, the observation marginal $p(o\mid b,a)=\int Z_\theta(o\mid s')P_\theta(s'\mid s,a)\,b(ds)\,ds'$, and define
\begin{equation}\label{eq:psi_tilde_f_def}
\psi_{o,a}(s):=\int Z_\theta(o\mid s')P_\theta(s'\mid s,a)\,ds', 
\qquad
\tilde f_{o,a}(s):=\int f(s')Z_\theta(o\mid s')P_\theta(s'\mid s,a)\,ds',
\end{equation}
so that $p(o\mid b,a)=b(\psi_{o,a})$ and $\tau(b,a,o)(f)=b(\tilde f_{o,a})/b(\psi_{o,a})$.

Under Assumption~\ref{aspt:particle_error}, for all $b,b'\in\Delta(\gS)$:
\begin{equation}\label{eq:Q_test_decomp}
|Q_\theta^{\star,B}(b,a)-Q_\theta^{\star,B}(b',a)|
\;\le\;
\sum_{\ell=0}^{H-1}\gamma^\ell\,C_{\mathrm{BU}}^\ell\sup_{f\in\gF_\ell(\theta,a)}|b(f)-b'(f)|
\;+\;2\gamma^H V_{\max},
\end{equation}
where $C_{\mathrm{BU}}:=2g_{\max}/g_{\min}^2$ and $\gF_\ell(\theta,a)$ is a class of functions $\gS\to[-g_{\max}^\ell V_{\max},\,g_{\max}^\ell V_{\max}]$ with $\gN_{[]}(\gF_\ell,\epsilon)\le|\gA|^\ell\gN_\Theta(\epsilon/(\ell V_{\max}))$, defined recursively:
\begin{itemize}
\item $\gF_0(\theta,a)=\{r(\cdot,a'):a'\in\gA\}\cup\{Q_\theta^{\star,S}(\cdot,a'):a'\in\gA\}$,
\item $\gF_{\ell+1}(\theta,a)=\bigl\{\tilde f_{o,a''}:\,f\in\gF_\ell(\theta,a''),\,a''\in\gA\bigr\}\cup\bigl\{\psi_{o,a''}:\,a''\in\gA\bigr\}$.
\end{itemize}
\end{lem}

\begin{proof}
We prove by induction on $h$: for any $\Phi(b)=\E_\theta^{\pi}[\sum_{u=0}^{h-1}\gamma^u r_u\mid b_0=b]$ with $h\le H$, $|\Phi(b)-\Phi(b')|$ admits the decomposition~\eqref{eq:Q_test_decomp}. Taking $\Phi=Q_\theta^{\star,B}(\cdot,a)$ with $h=H$ and bounding the tail by $2\gamma^H V_{\max}$ yields the result.

\emph{Base case ($h=1$).} $|\Phi(b)-\Phi(b')|=|\sup_a[b(r_a)-b'(r_a)]|\le\sup_{f\in\gF_0}|b(f)-b'(f)|$.

\emph{Inductive step ($h\to h+1$).} For a fixed action $a$, write
\[
\Phi(b)=b(r_a)+\gamma\int_{\gO}\Phi_{h-1}\bigl(\tau(b,a,o)\bigr)\,p(o\mid b,a)\,do,
\]
where $\Phi_{h-1}$ is the $(h-1)$-step continuation. The continuation difference decomposes as:
\begin{align}
&\left|\int\Phi_{h-1}(\tau(b,a,o))\,p(o|b,a)\,do-\int\Phi_{h-1}(\tau(b',a,o))\,p(o|b',a)\,do\right|\notag\\
&\;\le\;\underbrace{\|\Phi_{h-1}\|_\infty\int|(b-b')(\psi_{o,a})|\,do}_{\text{(I) obs.\ distribution shift}}+\underbrace{\int|\Phi_{h-1}(\tau(b,a,o))-\Phi_{h-1}(\tau(b',a,o))|\,p(o|b',a)\,do}_{\text{(II) belief propagation}},\label{eq:cont_split}
\end{align}
using $p(o\mid b,a)-p(o\mid b',a)=(b-b')(\psi_{o,a})$.

\emph{Term (I)} contributes $\psi_{o,a}\in\gF_{\ell+1}$. 

\emph{Term (II)} requires bounding $|\tau(b,a,o)(f)-\tau(b',a,o)(f)|$ for $\|f\|_\infty\le 1$. Since $\tau(b,a,o)(f)=b(\tilde f_{o,a})/b(\psi_{o,a})$, the quotient rule gives
\begin{equation}\label{eq:quotient_rule}
\frac{b(\tilde f_{o,a})}{b(\psi_{o,a})}-\frac{b'(\tilde f_{o,a})}{b'(\psi_{o,a})}
\;=\;
\frac{(b-b')(\tilde f_{o,a})}{b(\psi_{o,a})}-\frac{b'(\tilde f_{o,a})}{b(\psi_{o,a})\,b'(\psi_{o,a})}\,(b-b')(\psi_{o,a}).
\end{equation}
Since $b(\psi_{o,a}),b'(\psi_{o,a})\ge g_{\min}$ and $|b'(\tilde f_{o,a})|\le g_{\max}$:
\begin{equation}\label{eq:tau_test_fn}
|\tau(b,a,o)(f)-\tau(b',a,o)(f)|
\;\le\;
\frac{|(b-b')(\tilde f_{o,a})|}{g_{\min}}+\frac{g_{\max}}{g_{\min}^2}\,|(b-b')(\psi_{o,a})|.
\end{equation}
Both $\tilde f_{o,a}$ and $\psi_{o,a}$ are in $\gF_{\ell+1}$, with per-update amplification
\[
C_{\mathrm{BU}}=\frac{1}{g_{\min}}+\frac{g_{\max}}{g_{\min}^2}\;\le\;\frac{2\,g_{\max}}{g_{\min}^2},
\]
By the inductive hypothesis, $|\Phi_{h-1}(\tau(b,a,o))-\Phi_{h-1}(\tau(b',a,o))|$ decomposes into evaluations of $\tau(b,a,o)$ vs.\ $\tau(b',a,o)$ on $\gF_\ell$. Substituting~\eqref{eq:tau_test_fn} reduces each to evaluations of $b$ vs.\ $b'$ on $\gF_{\ell+1}$, with one factor of $C_{\mathrm{BU}}$ per level. Taking $\sup_a$ over the action set and iterating across $H$ levels yields~\eqref{eq:Q_test_decomp}.
\end{proof}

\begin{lem}[$Q^{\star,B}$ is Lipschitz in TV]\label{lem:Q_lipschitz_tv}
For any $\theta\in\Theta$, $a\in\gA$, and $b,b'\in\Delta(\gS)$,
\[
|Q_\theta^{\star,B}(b,a)-Q_\theta^{\star,B}(b',a)|\;\le\;2V_{\max}\,\|b-b'\|_{\mathrm{TV}}.
\]
\end{lem}

\begin{proof}
$|r(b,a)-r(b',a)|\le 2R_{\max}\|b-b'\|_{\mathrm{TV}}$. The trajectory law under $b$ and $b'$ differs in TV by at most $\|b-b'\|_{\mathrm{TV}}$ (data processing), so the $V_{\max}$-bounded continuation differs by at most $2V_{\max}\|b-b'\|_{\mathrm{TV}}$. Using $R_{\max}+\gamma V_{\max}=V_{\max}$ gives the result.
\end{proof}

\begin{lem}[Performance decomposition]\label{lem:perf_decomp_delta}
Fix $\theta\in\Theta$. Let $\pi_t$ denote the \emph{stationary} policy induced by the planner with frozen parameters (model $\bar\theta_{k(t)}$, belief $b_t^J$) at time $t$, and let $\delta_u:=V_{\theta}^{\star,B}(b_u)-Q_{\theta}^{\star,B}(b_u,A_u)$ where $A_u\sim\pi_t(\cdot\mid H_u)$. Then
\[
V_{\theta}^{\star,B}(b_t)-V_{\theta}^{\pi_t,B}(b_t)
\;\le\;
\E\!\left[\sum_{u=t}^\infty \gamma^{u-t}\delta_u\,\Bigm|\, H_t\right].
\]
\end{lem}

\begin{proof}
Unroll the Bellman recursion $V_{\theta}^{\star,B}(b_u)-V_{\theta}^{\pi_t,B}(b_u)=\delta_u+\gamma\,\E[V_{\theta}^{\star,B}(b_{u+1})-V_{\theta}^{\pi_t,B}(b_{u+1})\mid H_u]$ from $u=t$ to $u=t+n$, using $|V_{\theta}^{\star,B}-V_{\theta}^{\pi_t,B}|\le 2V_{\max}$ as the tail bound, and send $n\to\infty$.
\end{proof}

\begin{lem}[Per-step bound on $\delta_u$]\label{lem:delta_bound_main}
For every $u$ and any $d\ge 0$,
\[
\delta_u
\;\le\;
2\,d_Q(b_u,b_u^J)
+
2\,\epsilon_{\mathrm{mc}}(u,d)
+
2\,\epsilon_{\mathrm{struct}}(u,d).
\]
\end{lem}

\begin{proof}
Let $a_u^\star\in\arg\max_a Q_{\theta}^{\star,B}(b_u,a)$ and $\hat a_u\in\arg\max_a \widehat Q_u(b_u^J,a)$. Then
\begin{align*}
\delta_u
&=Q_{\theta}^{\star,B}(b_u,a_u^\star)-Q_{\theta}^{\star,B}(b_u,\hat a_u)\\
&\le \underbrace{|Q_{\theta}^{\star,B}(b_u,a_u^\star)-Q_{\theta}^{\star,B}(b_u^J,a_u^\star)|+|Q_{\theta}^{\star,B}(b_u^J,\hat a_u)-Q_{\theta}^{\star,B}(b_u,\hat a_u)|}_{\le\, 2\,d_Q(b_u,b_u^J) \;\text{ by Definition~\ref{def:dQ}}}\\
&\quad+Q_{\theta}^{\star,B}(b_u^J,a_u^\star)-Q_{\theta}^{\star,B}(b_u^J,\hat a_u).
\end{align*}
For the last line, add and subtract $Q_{\theta,d}$ and $\widehat Q_u$, using optimality of $\hat a_u$ for $\widehat Q_u$:
\[
Q_{\theta,d}(b_u^J,a_u^\star)
\le \widehat Q_u(b_u^J,a_u^\star)+\epsilon_{\mathrm{mc}}
\le \widehat Q_u(b_u^J,\hat a_u)+\epsilon_{\mathrm{mc}}
\le Q_{\theta,d}(b_u^J,\hat a_u)+2\epsilon_{\mathrm{mc}},
\]
so $Q_{\theta}^{\star,B}(b_u^J,a_u^\star)-Q_{\theta}^{\star,B}(b_u^J,\hat a_u)\le 2\epsilon_{\mathrm{struct}}+2\epsilon_{\mathrm{mc}}$.
\end{proof}

% \begin{rmk}[Belief dynamics convention]\label{rmk:belief_dynamics}
% In the decomposition~\eqref{eq:main_decomp_round_t}, $V_{\bar\theta_t}^{\star,B}(b_t)$ and $V_{\bar\theta_t}^{\pi_t,B}(b_t)$ are values in the belief-MDP \emph{under model $\bar\theta_t$}, evaluated at the true belief $b_t=b_t^{\theta^\star}$. The planning bound below unrolls the Bellman recursion under $\bar\theta_t$'s dynamics (both state transitions and belief updates). Term~(C) in~\eqref{eq:main_decomp_round_t} then measures the value difference of the \emph{same} policy $\pi_t$ between $\bar\theta_t$'s and $\theta^\star$'s belief-MDPs, which is exactly what the simulation lemma (Lemma~\ref{lem:simulation}) controls.
% \end{rmk}

\begin{thm}[Round-$t$ planning + state-identification bound]\label{thm:planning_state_bound}
For each $t\ge 1$ and any depth $d$,
\[
V_{\bar\theta_{k(t)}}^{\star,B}(b_t)-V_{\bar\theta_{k(t)}}^{\pi_t,B}(b_t)
\;\le\;
\E\!\left[\sum_{u=t}^\infty \gamma^{u-t}\bigl(2\,d_Q(b_u,b_u^J)
+2\,\epsilon_{\mathrm{mc}}(u,d)+2\,\epsilon_{\mathrm{struct}}(u,d)\bigr)\,\Bigm|\,H_t\right].
\]
\end{thm}

\begin{proof}
Substitute the per-step bound of Lemma~\ref{lem:delta_bound_main} into Lemma~\ref{lem:perf_decomp_delta}.
\end{proof}

\paragraph{Observation discretization.}
For continuous $\gO$, the tree never revisits observation-children (each sampled $o$ is unique), so every node below the root holds a single particle.
At $d=1$ each leaf value~\eqref{eq:qmdp_leaf_eval} then reduces to $\max_{a'}Q_\theta^{\star,S}(s',a')$ at a single state, which targets the linear surrogate rather than $Q_{\theta,1}$, and for $d\ge2$ the tree cannot build depth beyond~$1$.
We discretize: let $\phi:\gO\to\bar\gO$ with $|\bar\gO|=M$ and bin diameter $\mathrm{diam}(\phi):=\max_{\bar o}\sup_{o,o'\in\phi^{-1}(\bar o)}\|o-o'\|$.
Algorithms~\ref{alg:mctq} and~\ref{alg:simulate} are modified so that $\textsc{Child}(v,a',\phi(o'))$ replaces $\textsc{Child}(v,a',o')$.
The bin-averaged posterior is $\tau_\theta^{\phi}(b,a,\bar o):=\int_{\phi^{-1}(\bar o)} \tau_\theta(b,a,o)\,p_\theta(o\mid b,a)\,\mathrm do\big/\int_{\phi^{-1}(\bar o)} p_\theta(o\mid b,a)\,\mathrm do$, and we write $Q_{\theta,d}^{\phi}$ for the target~\eqref{eq:Qd_def} with $\tau_\theta(b,a,o)$ replaced by $\tau_\theta^{\phi}(b,a,\phi(o))$ at every level.
When $|\gO|<\infty$, take $\phi=\mathrm{id}$, $M=|\gO|$, $\mathrm{diam}(\phi)=0$.
When $d=0$, no tree is built and no discretization is needed.

\paragraph{Seeds and the estimator.}
Fix a round $t$ and a root action $a$, and write $\theta:=\bar\theta_{k(t)}$ and $b^J:=b_t^J$.
All probabilities and expectations below are conditional on $b_t^J$, and we take $d\ge1$.
Simulation $n\in[N]$ is driven by a \emph{seed}
\begin{equation}\label{eq:mc_seed}
\xi_n:=\bigl(j_n,\,u_{n,0},\,(a'_{n,\ell},u_{n,\ell})_{\ell=1}^{d-1}\bigr),
\end{equation}
where $j_n\sim\mathrm{Unif}([J])$ is the sampled particle index, $a'_{n,\ell}\sim\mathrm{Unif}(\gA)$ is the action drawn at depth $\ell$, and $u_{n,\ell}$ is the noise that generates the next state and observation at depth $\ell$ from $P_\theta$ and $Z_\theta$.
The seeds $\xi_1,\ldots,\xi_N$ are i.i.d.
Running Algorithm~\ref{alg:mctq} on them produces the tree $T_a$, and we write the root estimate as a deterministic function of all seeds,
\begin{equation}\label{eq:mc_fa}
\widehat Q(b^J,a)=f_a(\xi_1,\ldots,\xi_N):=\frac1N\sum_{n=1}^N\bigl[r(s_{j_n},a)+\gamma\,\hat V(v_{n,1})\bigr],
\end{equation}
where $v_{n,\ell}$ is the depth-$\ell$ node on the path of simulation $n$ and $\hat V$ is the DP backup~\eqref{eq:dp_backup}.
The function $f_a$ is not an average of independent terms, since $\hat V(v_{n,1})$ depends on every simulation that passes through $v_{n,1}$.
For $S\subseteq[N]$, let $f_a^{(-S)}$ be the same expression computed from the seeds $(\xi_m)_{m\notin S}$ alone, with the sum running over $n\notin S$ and the normalizer $1/N$ retained.

\paragraph{Balancedness.}
For $\rho\ge1$, let $\mathcal E_\rho$ be the event that every internal node $v$ of $T_a$ satisfies
\begin{equation}\label{eq:balanced}
N(v)\le\rho\,|\gA|\,N(v,a')\qquad\text{for all }a'\in\gA,
\end{equation}
so that every action receives at least a $1/(\rho|\gA|)$ fraction of the visits of every internal node.
Define the depth factor
\begin{equation}\label{eq:Lambda_d_def}
\Lambda_d:=\sum_{\ell=0}^{d-1}(\rho\gamma|\gA|)^{\ell},
\end{equation}
so that $\Lambda_0=0$ and $\Lambda_1=1$.
Let $p_\rho:=\max_{a\in\gA}\Pr(\mathcal E_\rho^c\mid b_t^J)$, where the event refers to the tree $T_a$, and let
\[
q_{\min}:=\inf_{s\in\gS,\,a''\in\gA,\,\bar o\in\bar\gO}\int_{\gS}Z_\theta\bigl(\phi^{-1}(\bar o)\mid s'\bigr)\,P_\theta(\mathrm ds'\mid s,a'')
\]
be the smallest one-step bin probability.
Under the bounded observation density of Assumption~\ref{aspt:particle_error}, $q_{\min}\ge g_{\min}\min_{\bar o}\lambda(\phi^{-1}(\bar o))$, where $\lambda$ is the reference measure of $Z_\theta$; for finite $\gO$ with the counting measure, $q_{\min}\ge g_{\min}$.

\begin{prop}[Monte Carlo error]\label{prop:mc_concentration}
Run Algorithm~\ref{alg:mctq} with uniform internal selection~\eqref{eq:uniform_action}, $N$ simulations per root action, observation discretization $\phi$ ($|\bar\gO|=M$), and returns in $[-V_{\max},V_{\max}]$, and let $\rho\ge1$.
Let
\begin{equation}\label{eq:bias_def}
\mathrm{bias}(N,d):=\sup_{a}\bigl|\E[\widehat Q(b^J,a)\mid b_t^J]-Q_{\bar\theta_{k(t)},d}(b^J,a)\bigr|
\end{equation}
and $B_d:=2V_{\max}\bigl(1+2(1+\gamma)\Lambda_d\bigr)$.
For every $\delta\in(0,1)$, with probability at least $1-\delta-|\gA|p_\rho$ conditionally on $b_t^J$,
\begin{equation}\label{eq:mc_hp}
\epsilon_{\mathrm{mc}}(t,d)
\;\le\;
\underbrace{2(1+\gamma)V_{\max}\Lambda_d\sqrt{\frac{2\log(2|\gA|/\delta)}{N}}}_{\text{fluctuation}}
\;+\;\mathrm{bias}(N,d)\;+\;B_d\,p_\rho,
\end{equation}
and, in second moment,
\begin{equation}\label{eq:mc_var_bias}
\E\bigl[\epsilon_{\mathrm{mc}}(t,d)^2\,\big|\,b_t^J\bigr]
\;\le\;
\frac{8(1+\gamma)^2\,|\gA|\,V_{\max}^2\,\Lambda_d^2}{N}
\;+\;
2\,|\gA|\,\mathrm{bias}(N,d)^2
\;+\;
3\,|\gA|\,B_d^2\,p_\rho.
\end{equation}
The bias satisfies
\begin{equation}\label{eq:bias_explicit}
\mathrm{bias}(N,d)
\;\le\;
\underbrace{\frac{c_0\,\Phi_d\,V_{\max}\sqrt{|\gA|}}{\sqrt{N}}}_{\text{DP backup estimation error}}
\;+\;
\underbrace{\frac{2\gamma(1-\gamma^d)}{1-\gamma}\cdot\frac{g_{\max}}{g_{\min}}\,V_{\max}\,\omega_Z(\mathrm{diam}(\phi))}_{\text{discretization error}},
\end{equation}
where $c_0>0$ is a constant independent of $N$, $M$ and $\delta$,
\begin{equation}\label{eq:Phi_d_def}
\Phi_d:=\sum_{\ell=1}^{d}(\gamma\sqrt{|\gA|M})^\ell,
\end{equation}
and $\omega_Z(\delta):=\sup_{\theta,s}\sup_{\|o-o'\|\le\delta}|Z_\theta(o\mid s)/Z_\theta(o'\mid s)-1|$.
If $d\le1$, then $p_\rho=0$.
If $d\ge2$, $q_{\min}>0$ and $|\gA|M\ge2$, then $p_2\le2|\gA|(|\gA|M)^d\exp\bigl(-N(q_{\min}/|\gA|)^{d-1}/32\bigr)$.
When $d=0$, $\widehat Q(b^J,a)=\frac1J\sum_{j}Q_\theta^{\star,S}(s_j,a)$ is exact and every term vanishes ($\Lambda_0=\Phi_0=0$).
When $\phi=\mathrm{id}$ ($M=|\gO|$, $\mathrm{diam}(\phi)=0$), the discretization error vanishes.
\end{prop}

No stationarity assumption on the rewards at internal nodes is required, and the simulation returns need not be independent.
The proof uses only that the seeds are i.i.d., that the path of each simulation depends on its own seed alone, and that the estimator has bounded differences on $\mathcal E_\rho$.
It rests on the following four results; throughout, we fix a root action $a$ and $d\ge1$.

\begin{lem}[Path decoupling]\label{lem:path_decoupling}
There is a deterministic map $\nu$ such that, whatever the other seeds are, the path of simulation $m$ (its states, actions, binned observations, and nodes $v_{m,1},\ldots,v_{m,d}$) equals $\nu(\xi_m)$.
\end{lem}

\begin{proof}
We induct on depth.
Simulation $m$ starts from $s_{j_m}$ with the root action $a$, and $u_{m,0}$ generates its next state and observation, so $v_{m,1}=(a,\phi(o))$ is a function of $\xi_m$.
Suppose the state $s_{m,\ell}$ and the node $v_{m,\ell}$ at depth $\ell<d$ are functions of $\xi_m$.
Whether or not an earlier simulation expanded $v_{m,\ell}$, Algorithm~\ref{alg:simulate} only initializes statistics and continues, and it selects the action $a'_{m,\ell}$, a component of $\xi_m$ that reads no tree statistic.
The next state and observation are generated from $(s_{m,\ell},a'_{m,\ell},u_{m,\ell})$, and $v_{m,\ell+1}$ is the history of $v_{m,\ell}$ extended by $a'_{m,\ell}$ and the binned observation.
Hence $s_{m,\ell+1}$ and $v_{m,\ell+1}$ are functions of $\xi_m$, and the simulation stops at depth $d$ regardless of the other seeds.
\end{proof}

Both ingredients of the lemma are needed.
Under UCB1, the action at $v_{m,\ell}$ depends on the statistics left by simulations $1,\ldots,m-1$.
Under a POMCP-style early return (``if $v$ is not yet expanded, expand it and return''), simulation $m$ stops at the first node that no earlier simulation created, so its depth depends on $\xi_1,\ldots,\xi_{m-1}$.

\begin{cor}[Localized influence]\label{cor:two_paths}
Let $S\subseteq[N]$ and let $P_S$ be the set of nodes on the paths $\nu(\xi_m)$, $m\in S$.
Every node $v\notin P_S$ has the same particle set $B(v)$, counts $N(v)$ and $N(v,a')$, mean rewards $\bar r(v,a')$, and DP value $\hat V(v)$ in the runs on $(\xi_m)_{m\in[N]}$ and on $(\xi_m)_{m\notin S}$.
In particular, replacing one seed $\xi_n$ by $\xi_n'$ changes these quantities only at nodes on the two root-to-leaf paths $\nu(\xi_n)$ and $\nu(\xi_n')$.
\end{cor}

\begin{proof}
By Lemma~\ref{lem:path_decoupling}, in either run the simulations through $v$ are those $m$ with $v\in\nu(\xi_m)$ among the seeds present.
For $v\notin P_S$ no $m\in S$ is among them, so $B(v)$, $N(v)$, $N(v,a')$ and $\bar r(v,a')$ coincide.
A descendant of $v$ lies on a path in $P_S$ only if $v$ does, so the whole subtree of $v$ avoids $P_S$, and backward induction on depth with~\eqref{eq:dp_backup} gives the same $\hat V(v)$.
For the last claim, compare both runs with the run on $(\xi_m)_{m\ne n}$.
\end{proof}

\begin{lem}[Bounded difference]\label{lem:bounded_difference}
On $\mathcal E_\rho$, for every $S\subseteq[N]$,
\begin{equation}\label{eq:bd_removal}
\bigl|f_a(\xi_1,\ldots,\xi_N)-f_a^{(-S)}\bigr|\le|S|\cdot\frac{2(1+\gamma)V_{\max}\Lambda_d}{N}.
\end{equation}
Consequently, if two seed vectors $\xi,\xi'$ differ in $k$ coordinates and both of their runs lie in $\mathcal E_\rho$, then $|f_a(\xi)-f_a(\xi')|\le k\,c$, where
\begin{equation}\label{eq:bd_constant}
c:=\frac{4(1+\gamma)V_{\max}\Lambda_d}{N}.
\end{equation}
\end{lem}

\begin{proof}
Write $K(v):=N(v)$ in the full run, $k(v)$ for the number of simulations $m\in S$ whose path passes through $v$, and $\Delta(v):=|\hat V(v)-\hat V^{(-S)}(v)|$, where $\hat V^{(-S)}$ is the DP value in the reduced run.
If $k(v)=K(v)$, node $v$ is absent from the reduced run, and we set $\hat V^{(-S)}(v):=0$.
By Corollary~\ref{cor:two_paths}, $\Delta(v)=0$ unless $k(v)\ge1$.
Define the weighted influence $I(v):=(K(v)-k(v))\,\Delta(v)$ and the sequence
\[
y_d:=1,\qquad y_\ell:=\rho|\gA|\,(1+\gamma y_{\ell+1})\quad(1\le\ell\le d-1).
\]
We show by backward induction on depth that every node $v$ at depth $\ell\in[1,d]$ satisfies
\begin{equation}\label{eq:influence_induction}
I(v)\le 2V_{\max}\,y_\ell\,k(v).
\end{equation}

\emph{Leaves ($\ell=d$).}
If $k(v)=K(v)$, then $I(v)=0$.
Otherwise, for each $a'$, removing $k(v)$ of the $K(v)$ states in $B(v)$ moves the average of $Q_\theta^{\star,S}(\cdot,a')$ by
\[
\frac{k(v)}{K(v)}\,\bigl|\text{mean over removed states}-\text{mean over remaining states}\bigr|\le\frac{2V_{\max}\,k(v)}{K(v)}.
\]
The maximum over $a'$ in~\eqref{eq:qmdp_leaf_eval} is $1$-Lipschitz in the sup norm, so $\Delta(v)\le 2V_{\max}k(v)/K(v)$ and $I(v)\le 2V_{\max}k(v)$.

\emph{Internal nodes ($1\le\ell\le d-1$).}
Fix $a'$, and write $K':=N(v,a')$ and $k'$ for the number of simulations in $S$ that select $a'$ at $v$.
If $k'=0$, then $\hat Q(v,a')$ is unchanged, since no child of $v$ under $a'$ carries a simulation of $S$.
If $1\le k'<K'$, let $\Sigma$ and $\Sigma'$ be the sums of $r(s,a')+\gamma\hat V(\text{child})$ over the simulations that select $a'$ at $v$ in the full and in the reduced run.
Each removed simulation contributes a term of absolute value at most $V_{\max}$, because $|r|\le(1-\gamma)V_{\max}$ and $|\hat V|\le V_{\max}$.
Each of the $K(w)-k(w)$ remaining simulations entering a child $w$ changes by $\gamma\Delta(w)$.
Hence
\[
|\Sigma-\Sigma'|\le k'V_{\max}+\gamma\sum_{w\in\mathrm{Ch}(v,a')}I(w).
\]
Since $\hat Q(v,a')=\Sigma/K'$, the reduced-run value is $\Sigma'/(K'-k')$, and $|\Sigma'|\le(K'-k')V_{\max}$,
\[
\Bigl|\frac{\Sigma}{K'}-\frac{\Sigma'}{K'-k'}\Bigr|
\le\frac{|\Sigma-\Sigma'|}{K'}+\frac{k'\,|\Sigma'|}{K'(K'-k')}
\le\frac{1}{K'}\Bigl(2k'V_{\max}+\gamma\sum_{w\in\mathrm{Ch}(v,a')}I(w)\Bigr).
\]
If $k'=K'$, action $a'$ is unvisited in the reduced run and drops out of the maximum in~\eqref{eq:dp_backup}.
A maximum of values in $[-V_{\max},V_{\max}]$ changes by at most the largest change of its retained arguments, or by at most $2V_{\max}=2k'V_{\max}/K'$ if its maximizer drops out.
Writing $k'_{a'}$ for the $k'$ of action $a'$, in every case
\[
\Delta(v)\le\max_{a'\in\gA}\frac{1}{N(v,a')}\Bigl(2k'_{a'}V_{\max}+\gamma\sum_{w\in\mathrm{Ch}(v,a')}I(w)\Bigr).
\]
Multiplying by $K(v)$, using~\eqref{eq:balanced}, and bounding the maximum by the sum,
\[
I(v)\le K(v)\,\Delta(v)\le\rho|\gA|\sum_{a'\in\gA}\Bigl(2k'_{a'}V_{\max}+\gamma\sum_{w\in\mathrm{Ch}(v,a')}I(w)\Bigr).
\]
By the induction hypothesis~\eqref{eq:influence_induction} at depth $\ell+1$ and $\sum_{a'}k'_{a'}=\sum_{a'}\sum_{w\in\mathrm{Ch}(v,a')}k(w)=k(v)$,
\[
I(v)\le\rho|\gA|\,\bigl(2V_{\max}+2\gamma V_{\max}\,y_{\ell+1}\bigr)\,k(v)=2V_{\max}\,y_\ell\,k(v).
\]

\emph{Root.}
Each simulation enters exactly one depth-$1$ node, and the normalizer $1/N$ is the same in both runs, so
\[
\bigl|f_a-f_a^{(-S)}\bigr|
\le\frac1N\sum_{m\in S}\bigl|r(s_{j_m},a)+\gamma\hat V(v_{m,1})\bigr|+\frac{\gamma}{N}\sum_{w:\,\text{depth }1}I(w)
\le\frac{|S|}{N}\,V_{\max}\,(1+2\gamma y_1).
\]
Unrolling the recursion,
\[
y_1=\rho|\gA|\sum_{i=0}^{d-2}(\rho\gamma|\gA|)^i+(\rho\gamma|\gA|)^{d-1}.
\]
By~\eqref{eq:Lambda_d_def} and $(\rho\gamma|\gA|)^{d-1}\le\Lambda_d$,
\[
1+2\gamma y_1=2\Lambda_d-1+2\gamma(\rho\gamma|\gA|)^{d-1}\le2(1+\gamma)\Lambda_d.
\]
This proves~\eqref{eq:bd_removal}.
For the second claim, let $S$ be the set of coordinates where $\xi$ and $\xi'$ differ.
Both runs have the same reduced run on $(\xi_m)_{m\notin S}$, so~\eqref{eq:bd_removal} applied to each and the triangle inequality give $|f_a(\xi)-f_a(\xi')|\le 2k\cdot2(1+\gamma)V_{\max}\Lambda_d/N=kc$.
\end{proof}

\begin{lem}[Fluctuation up to a bad event]\label{lem:mc_fluctuation}
Let $p:=\Pr(\mathcal E_\rho^c\mid b_t^J)$ for the tree $T_a$, and let $B_d:=2V_{\max}+cN$ with $c$ from~\eqref{eq:bd_constant}.
For every $\delta\in(0,1)$, with probability at least $1-\delta-p$,
\[
\bigl|f_a-\E f_a\bigr|\le 2(1+\gamma)V_{\max}\Lambda_d\sqrt{\frac{2\log(2/\delta)}{N}}+B_d\,p,
\]
and for every $q\in[-V_{\max},V_{\max}]$,
\[
\E\bigl[(f_a-q)^2\bigr]\le\frac{8(1+\gamma)^2V_{\max}^2\Lambda_d^2}{N}+2\,(\E f_a-q)^2+3B_d^2\,p.
\]
\end{lem}

\begin{proof}
Let $G$ be the set of seed vectors whose run lies in $\mathcal E_\rho$.
If $\Pr(G)=0$, both claims are trivial, so assume otherwise.
Let $H$ denote the Hamming distance and define the extension
\[
g(\xi):=\inf_{\eta\in G}\bigl\{f_a(\eta)+c\,H(\xi,\eta)\bigr\}.
\]
First, $g=f_a$ on $G$: for $\xi\in G$, Lemma~\ref{lem:bounded_difference} gives $f_a(\eta)+cH(\xi,\eta)\ge f_a(\xi)$ for every $\eta\in G$, with equality at $\eta=\xi$.
Second, $g$ changes by at most $c$ when one coordinate of $\xi$ changes, because each $H(\cdot,\eta)$ does.
Third, $-V_{\max}\le g\le V_{\max}+cN$, so $|f_a-g|\le B_d$ everywhere and
\[
|\E f_a-\E g|\le\E\bigl[|f_a-g|\,\mathbf{1}_{G^c}\bigr]\le B_d\,p.
\]
The seeds are independent, so McDiarmid's bounded-differences inequality applies to $g$:
\[
\Pr\bigl(|g-\E g|\ge s\bigr)\le2\exp\Bigl(-\frac{2s^2}{Nc^2}\Bigr).
\]
Setting the right-hand side to $\delta$ gives $s=c\sqrt{N\log(2/\delta)/2}=2(1+\gamma)V_{\max}\Lambda_d\sqrt{2\log(2/\delta)/N}$.
On $G\cap\{|g-\E g|<s\}$, an event of probability at least $1-\delta-p$,
\[
|f_a-\E f_a|\le|g-\E g|+|\E g-\E f_a|<s+B_d\,p.
\]
For the second moment, the Efron--Stein inequality gives $\mathrm{Var}(g)\le Nc^2/2=8(1+\gamma)^2V_{\max}^2\Lambda_d^2/N$.
Since $f_a=g$ on $G$ and $|f_a-q|\le2V_{\max}$,
\[
\E\bigl[(f_a-q)^2\bigr]\le\E\bigl[(g-q)^2\bigr]+4V_{\max}^2\,p=\mathrm{Var}(g)+(\E g-q)^2+4V_{\max}^2\,p.
\]
Finally, $(\E g-q)^2\le2(\E f_a-q)^2+2B_d^2p^2$, and $2B_d^2p^2+4V_{\max}^2p\le3B_d^2p$.
The infimum defining $g$ is universally measurable, which suffices for both inequalities.
\end{proof}

\begin{lem}[Balancedness under uniform selection]\label{lem:balance}
If $d\le1$, then $\mathcal E_\rho$ holds surely.
If $d\ge2$, $q_{\min}>0$ and $|\gA|M\ge2$, then
\[
\Pr(\mathcal E_2^c\mid b_t^J)\le2|\gA|\,(|\gA|M)^d\exp\Bigl(-\frac{N}{32}\Bigl(\frac{q_{\min}}{|\gA|}\Bigr)^{d-1}\Bigr).
\]
\end{lem}

\begin{proof}
For $d\le1$ the tree has no internal node.
Let $d\ge2$ and fix a potential internal node $v$ of $T_a$ at depth $\ell\in[1,d-1]$, i.e., a history with $\ell$ binned observations and $\ell-1$ internal actions.
By Lemma~\ref{lem:path_decoupling}, the indicators that simulation $m$ reaches $v$, and that it reaches $v$ and selects $a'$ there, are functions of $\xi_m$.
They are therefore i.i.d.\ Bernoulli across $m$, with parameters $p_v$ and $p_v/|\gA|$.
Each binned observation along the history has conditional probability at least $q_{\min}$ and each internal action has probability $1/|\gA|$, so
\[
\frac{p_v}{|\gA|}\ge\Bigl(\frac{q_{\min}}{|\gA|}\Bigr)^{\ell}\ge\pi:=\Bigl(\frac{q_{\min}}{|\gA|}\Bigr)^{d-1}.
\]
The multiplicative Chernoff bounds give
\begin{align*}
\Pr\bigl(N(v)\ge\tfrac32Np_v\bigr)&\le e^{-Np_v/12}\le e^{-N\pi/32},\\
\Pr\bigl(N(v,a')\le\tfrac34Np_v/|\gA|\bigr)&\le e^{-Np_v/(32|\gA|)}\le e^{-N\pi/32}.
\end{align*}
Outside these events, $N(v)<\frac32Np_v=2|\gA|\cdot\frac34Np_v/|\gA|<2|\gA|\,N(v,a')$, which is~\eqref{eq:balanced} with $\rho=2$.
A union bound over the $|\gA|$ actions and over the at most $\sum_{\ell=1}^{d-1}M^\ell|\gA|^{\ell-1}\le(|\gA|M)^d$ potential internal nodes gives the claim, since $1+|\gA|\le2|\gA|$.
\end{proof}

\begin{proof}[Proof of Proposition~\ref{prop:mc_concentration}]
When $d=0$ the planner is exact, so let $d\ge1$ and write $q_a:=Q_{\theta,d}(b^J,a)\in[-V_{\max},V_{\max}]$.

\emph{Step 1: Fluctuation.}
Apply Lemma~\ref{lem:mc_fluctuation} to each root action with $\delta/|\gA|$ in place of $\delta$, and take a union bound over $a\in\gA$.
With probability at least $1-\delta-|\gA|p_\rho$, every $a$ satisfies
\begin{align*}
|\widehat Q(b^J,a)-q_a|
&\le|\widehat Q(b^J,a)-\E\widehat Q(b^J,a)|+|\E\widehat Q(b^J,a)-q_a|\\
&\le2(1+\gamma)V_{\max}\Lambda_d\sqrt{\frac{2\log(2|\gA|/\delta)}{N}}+B_d\,p_\rho+\mathrm{bias}(N,d),
\end{align*}
and taking the supremum over $a$ gives~\eqref{eq:mc_hp}.
For the second moment, Lemma~\ref{lem:mc_fluctuation} with $q=q_a$ gives
\[
\E\bigl[(\widehat Q(b^J,a)-q_a)^2\bigr]\le\frac{8(1+\gamma)^2V_{\max}^2\Lambda_d^2}{N}+2\,\mathrm{bias}(N,d)^2+3B_d^2\,p_\rho.
\]
Summing over $a\in\gA$ through $\epsilon_{\mathrm{mc}}(t,d)^2\le\sum_a(\widehat Q(b^J,a)-q_a)^2$ gives~\eqref{eq:mc_var_bias}.
The bounds on $p_\rho$ are Lemma~\ref{lem:balance}.

\emph{Step 2: DP backup estimation error.}
The simulations follow the uniform policy at internal nodes, and the DP backup $\hat V(v)=\max_{a'}\hat Q(v,a')$ recovers the optimal internal actions.
We bound the bias $\mathrm{bias}_{\mathrm{DP}}(N,d):=\sup_a|\E\widehat Q(b^J,a)-Q_{\theta,d}^{\phi}(b^J,a)|$ of the resulting estimate by backward induction on depth, at the level of orders of magnitude, without any independence assumption on the returns.
Fix a node $v$ at depth $\ell\in[1,d]$ with $n_v:=N(v)$ visits, and condition on the set of simulations that reach $v$ and on their states there.
By Lemma~\ref{lem:path_decoupling}, the continuation of each such simulation below $v$ is driven by the components of its own seed at depths $\ge\ell$, which are independent across simulations and independent of the conditioning.
Conditionally, the states in $B(v)$ are therefore independent draws from the binned belief $b_v$ at $v$, and each $\hat Q(v,a')$ is formed from $n_{v,a'}:=N(v,a')$ independent continuations, with $n_{v,a'}\ge n_v/(\rho|\gA|)$ on $\mathcal E_\rho$.
Concentration of these conditional averages gives
\[
\bigl|\hat Q(v,a')-Q_{\theta,d-\ell}^{\phi}(b_v,a')\bigr|
\;\le\;
O\!\left(\frac{V_{\max}}{\sqrt{n_{v,a'}}}\right)
\;\le\;
O\!\left(\frac{V_{\max}\sqrt{|\gA|}}{\sqrt{n_v}}\right).
\]
At a leaf ($\ell=d$) the same bound holds for the pooled averages $\frac{1}{|B(v)|}\sum_{s\in B(v)}Q_\theta^{\star,S}(s,a')$ around $\int Q_\theta^{\star,S}(s,a')\,b_v(\mathrm ds)$.
This level must be included, because the maximum over $a'$ in~\eqref{eq:qmdp_leaf_eval} is biased upward at every finite $|B(v)|$.
Under the branching structure, a depth-$\ell$ node receives $n_v=\Omega(N/(|\gA|M)^\ell)$ visits in expectation when the bins have comparable probabilities.
Taking the maximum over estimated Q-values, the backup error at depth $\ell$ contributes at most $O(V_{\max}\sqrt{|\gA|^{\ell+1}M^\ell/N})$, discounted by $\gamma^\ell$.
The complement of $\mathcal E_\rho$ changes the conditional mean by at most $2V_{\max}p_\rho$, which is absorbed into the constant once $p_\rho\le N^{-1/2}$ (Remark~\ref{rmk:mc_constants}).
Summing over depths, leaves included,
\begin{equation}\label{eq:bias_derivation}
\mathrm{bias}_{\mathrm{DP}}(N,d)
\;\le\;\sum_{\ell=1}^{d}\gamma^\ell\cdot\frac{c_1 V_{\max}\sqrt{|\gA|}}{\sqrt{N/(|\gA|M)^\ell}}
\;=\;\frac{c_0\,\Phi_d\,V_{\max}\sqrt{|\gA|}}{\sqrt{N}}.
\end{equation}

\emph{Step 3: Discretization error.}
For $o,o'$ in the same bin, the Bayes update sensitivity~\citep[Section~4.3]{doucmoulinesolsson2014} gives
\[
\|\tau_\theta(b,a,o)-\tau_\theta(b,a,o')\|_{\mathrm{TV}}
\;\le\;
\frac{g_{\max}}{g_{\min}}\,\omega_Z(\mathrm{diam}(\phi)).
\]
Since $\tau_\theta^{\phi}$ is a convex combination of $\{\tau_\theta(b,a,o')\}_{o'\in\phi^{-1}(\bar o)}$, the same bound holds for $\|\tau_\theta^{\phi}(b,a,\bar o)-\tau_\theta(b,a,o)\|_{\mathrm{TV}}$.
By Lipschitz continuity of $\bar V_{\theta,d-1}^{\star,S}$ in TV (constant $2V_{\max}$, cf.\ Lemma~\ref{lem:Q_lipschitz_tv}), each of the $d$ tree levels contributes $\frac{2g_{\max}}{g_{\min}}V_{\max}\,\omega_Z(\mathrm{diam}(\phi))$ discounted by $\gamma$:
\begin{equation}\label{eq:disc_error_total}
\sup_{b,a}\bigl|Q_{\theta,d}^{\phi}(b,a)-Q_{\theta,d}(b,a)\bigr|
\;\le\;
\frac{2\gamma(1-\gamma^d)}{1-\gamma}\cdot\frac{g_{\max}}{g_{\min}}\,V_{\max}\,\omega_Z(\mathrm{diam}(\phi)).
\end{equation}
Combining Steps~2 and~3 gives $\mathrm{bias}(N,d)\le\mathrm{bias}_{\mathrm{DP}}(N,d)+\mathrm{bias}_{\mathrm{disc}}(d)$, yielding~\eqref{eq:bias_explicit}.
\end{proof}

\begin{rmk}[Constants, a verification-free variant, and the choice of depth]\label{rmk:mc_constants}
\emph{(a) Order of the constants.}
Compared with an analysis that treats the simulation returns as i.i.d., the only changes are the depth factor $(1+\gamma)\Lambda_d$ in the fluctuation term and the exponentially small remainder in $p_\rho$.
Since $\Lambda_d\le d\max(1,\rho\gamma|\gA|)^{d-1}$ and $\Phi_d\ge(\gamma\sqrt{|\gA|M})^d$, the fluctuation term of~\eqref{eq:mc_var_bias} is at most of the same exponential-in-$d$ order as the DP-backup term whenever $M=\Omega(\rho^2|\gA|)$.
Taking square roots in~\eqref{eq:mc_var_bias}, with $\mathrm{diam}(\phi)$ chosen as in Remark~\ref{rmk:disc_resolution}, then gives $\E[\epsilon_{\mathrm{mc}}(t,d)\mid b_t^J]\le C_{\mathrm{mc}}V_{\max}/\sqrt N$ with $C_{\mathrm{mc}}$ of the order stated in Theorem~\ref{thm:final_regret_bound}, whose order is therefore unchanged.
For coarse binning, $M\ll|\gA|$, the fluctuation dominates and $C_{\mathrm{mc}}\propto\sqrt{|\gA|}\,\Lambda_d$.
By Lemma~\ref{lem:balance}, the remainder $3|\gA|B_d^2p_2$ is at most $V_{\max}^2/N$ as soon as $N\ge32(|\gA|/q_{\min})^{d-1}\log\bigl(24N|\gA|^2(|\gA|M)^d(1+4\Lambda_d)^2\bigr)$, and it vanishes for $d\le1$.
At $d=0$, $\Lambda_0=0$ and the fluctuation vanishes identically.
For an internal rule satisfying (S1)--(S2) in place of uniform selection, the same proof applies with $|\gA|$ replaced by $1/\pi_{\min}$ in~\eqref{eq:balanced}, \eqref{eq:Lambda_d_def} and Lemma~\ref{lem:balance}.

\emph{(b) A verification-free variant.}
Sample splitting restores exactly conditionally i.i.d.\ returns.
Use $N/2$ seeds to build the tree, populate the particle sets, and run the DP backup, then freeze the leaf values and the greedy internal actions $\arg\max_{a'}\hat Q(v,a')$.
Value each root action with $N/2$ fresh seeds that follow the frozen greedy actions and stop at the frozen leaf values; a simulation that leaves the frozen tree stops and is evaluated by $\max_{a'}Q_\theta^{\star,S}(s,a')$ at its current state $s$.
Conditionally on the first half, the second-half returns are i.i.d.\ and bounded by $V_{\max}$, so Hoeffding's inequality applies to them verbatim at the cost of a factor $2$ in $N$, and the bias analysis of Step~2 applies to the first half.
Evaluating each leaf at its single incoming state instead of~\eqref{eq:qmdp_leaf_eval} removes the pooling but not the coupling through the maximum in the DP backup, so it restores independent returns only at $d=1$, where it also changes the leaf target to the linear surrogate $\widetilde V_\theta^{\star}$.

\emph{(c) Choice of depth.}
Propositions~\ref{prop:structural_error} and~\ref{prop:mc_concentration} concern different objects.
The former bounds the planning target, $\|\bar V_{\theta,d}^{\star,S}-V_\theta^{\star,B}\|_\infty\le\gamma^{d+1}\Xi_\theta/(1-\gamma)$, which involves no particles and decays geometrically in $d$.
The latter bounds the estimator $\widehat Q$ around $Q_{\theta,d}$, and its fluctuation $2(1+\gamma)V_{\max}\Lambda_d\sqrt{2\log(2|\gA|/\delta)/N}$ and DP-backup bias $c_0\Phi_dV_{\max}\sqrt{|\gA|/N}$ grow geometrically in $d$.
Their sum is U-shaped in $d$: depth buys structural bias at the price of Monte Carlo error, so each budget $N$ has a finite best depth.
\end{rmk}

\begin{rmk}[Optimal discretization resolution]\label{rmk:disc_resolution}
The discretization error in~\eqref{eq:bias_explicit} is a fixed bias $O(\omega_Z(\mathrm{diam}(\phi)))$, while the DP backup estimation error vanishes as $N\to\infty$ at rate $\Phi_d/\sqrt{N}$.
To make the discretization bias $O(1/\sqrt{N})$, set $\mathrm{diam}(\phi)=\omega_Z^{-1}\bigl((1-\gamma)g_{\min}/(2\gamma g_{\max}V_{\max}\sqrt{N})\bigr)$.
When $Z_\theta(o\mid s)$ is $L_Z$-Lipschitz in $o$ uniformly (so $\omega_Z(\delta)\le L_Z\delta/g_{\min}$) and $\gO\subseteq\R^{d_o}$ with a uniform grid:
\begin{equation}\label{eq:M_opt}
M=\left\lceil\left(\frac{2\gamma\,L_Z\,g_{\max}\,V_{\max}\,\mathrm{diam}(\gO)\sqrt{N}}{(1-\gamma)\,g_{\min}^2}\right)^{d_o}\right\rceil.
\end{equation}
With this choice, $M$ grows as $N^{d_o/2}$, and the simulation budget threshold for the DP backup error to be $O(1/\sqrt{N})$ is $N\ge c_0^2|\gA|\Phi_d^2$, which is exponential in $d\cdot d_o$ via the dependence of $\Phi_d$ on $M$.

Balancing the two bias terms instead gives an explicit rate.
Take $\gO\subseteq[0,D]^{d_o}$ with a cubic grid of mesh $h$, so that $M=(D/h)^{d_o}$ and $\mathrm{diam}(\phi)=\sqrt{d_o}\,h$, and assume $\omega_Z(\delta)\le L_Z\delta/g_{\min}$ as above.
For $\gamma\sqrt{|\gA|M}\ge2$ we have $\Phi_d\le2(\gamma\sqrt{|\gA|M})^d$, so the DP-backup and discretization terms of~\eqref{eq:bias_explicit} are of order $h^{-d\,d_o/2}N^{-1/2}$ and $h$, respectively, with constants depending on $d$, $\gamma$, $|\gA|$, $D$, $L_Z$, $g_{\max}/g_{\min}$ and $V_{\max}$.
Balancing them gives
\begin{equation}\label{eq:h_opt}
h^\star\asymp N^{-1/(d\,d_o+2)},\qquad \epsilon_{\mathrm{mc}}(t,d)=\widetilde O\bigl(N^{-1/(d\,d_o+2)}\bigr),
\end{equation}
since the fluctuation term of~\eqref{eq:mc_hp} does not depend on $h$ and is of lower order.
The $N^{-1/2}$ rate is recovered exactly when $d\,d_o=0$, i.e., at $d=0$.
At $d=1$, \eqref{eq:h_opt} is the familiar rate $N^{-1/(d_o+2)}$ of a histogram estimate, here of the leaf value as a function of the observation.
The balancedness condition of Lemma~\ref{lem:balance} holds at this mesh for large $N$: with the Lebesgue reference measure $q_{\min}\ge g_{\min}h^{d_o}$, so $N(q_{\min}/|\gA|)^{d-1}=\Omega(N^{(d_o+2)/(d\,d_o+2)})\to\infty$.
The computational cost is $N|\gA|$ simulations over at most $(|\gA|M)^d$ nodes.

When $d=0$ or $|\gO|<\infty$, this issue does not arise: set $\phi=\mathrm{id}$ and $M=|\gO|$.
\end{rmk}



%% ===================================================================
\subsection{Particle Filter Stability}
\label{app:pf_proof}

We work under Assumption~\ref{aspt:particle_error} throughout. Write $b_t:=b_t^{\theta^\star}$ for the true posterior and $b_t^J:=b_t^{J,(\bar\theta_{k(s)})_{s\le t}}$ for the particle belief filtered under the epoch-frozen model sequence. Since the model sequence is determined by the algorithm, we abbreviate the latter as $b_t^J$.

    \begin{prop}[Particle filter stability in $d_Q$]\label{prop:particle_error}
    Under Assumption~\ref{aspt:particle_error}, there exist constants $C_{\mathrm{pf}},C_{\mathrm{filt}}<\infty$ and a forgetting rate $\rho:=1-\varepsilon_{\mathrm{mix}}\,(g_{\min}/g_{\max})\in(0,1)$ such that for all $t\ge 1$,
    \begin{equation}\label{eq:pf_dQ_main}
    \E\bigl[d_Q(b_t,b_t^J)\bigr]
    \;\le\;
    \underbrace{C_{\mathrm{pf}}\cdot J^{-\frac{1}{2\kappa_{\mathrm{pf}}}}}_{\text{particle approx.}}
    \;+\;
    \underbrace{2V_{\max}\,C_{\mathrm{filt}}\sum_{s=1}^t \rho^{t-s}\,d_{\mathrm{model}}^{a_s}(\bar\theta_s,\theta^\star)}_{\text{filter mismatch}},
    \end{equation}
    where $d_{\mathrm{model}}^{a}(\theta,\theta'):=\sup_{s'}\|P_\theta(\cdot\mid s',a)-P_{\theta'}(\cdot\mid s',a)\|_{\mathrm{TV}}+\sup_{s',o}|Z_\theta(o\mid s')/Z_{\theta'}(o\mid s')-1|$, $C_{\mathrm{filt}}:=g_{\max}/g_{\min}$, and $\kappa_{\mathrm{pf}}$ is the \emph{particle filter exponent}:
    \begin{equation}\label{eq:kappa_pf_def}
    \kappa_{\mathrm{pf}}:=\max\!\left(1,\;\frac{\log(2g_{\max}^2/g_{\min}^2)}{\log(1/\gamma)}\right).
    \end{equation}
    The truncation depth $H_{\mathrm{pf}}$ in Lemma~\ref{lem:Q_lipschitz} is chosen to balance the sum and tail:
    \begin{equation}\label{eq:Hpf_choice}
    H_{\mathrm{pf}}:=\left\lceil\frac{\log(4V_{\max}\sqrt{J})}{\kappa_{\mathrm{pf}}\log(1/\gamma)}\right\rceil
    =\left\lceil\frac{\log(4V_{\max}\sqrt{J})}{\log\max(1/\gamma,\;2g_{\max}^2/g_{\min}^2)}\right\rceil,
    \end{equation}
    and $C_{\mathrm{pf}}$ depends polynomially on $V_{\max}$, $g_{\max}/g_{\min}$, $1/\varepsilon_{\mathrm{mix}}$, $\log\gN_\Theta(1/T)$, and $H_{\mathrm{pf}}$. When $\kappa_{\mathrm{pf}}=1$ (i.e., $2\gamma g_{\max}^2/g_{\min}^2\le 1$), the rate recovers $J^{-1/2}$. This $H_{\mathrm{pf}}$ is distinct from the trajectory truncation horizon $H$ in Section~\ref{app:model_error_proofs}.
    \end{prop}

    \begin{proof}[Proof of Proposition~\ref{prop:particle_error}]
    By the triangle inequality and Lemma~\ref{lem:Q_lipschitz_tv}:
    \begin{equation}\label{eq:pf_triangle}  
    d_Q(b_t^{\theta^\star},\,b_t^{J,\bar\theta_{1:t}})
    \;\le\;
    \underbrace{d_Q(b_t^{\theta^\star},\,b_t^{\bar\theta_{1:t}})}_{\text{(I) filter mismatch}}
    +
    \underbrace{d_Q(b_t^{\bar\theta_{1:t}},\,b_t^{J,\bar\theta_{1:t}})}_{\text{(II) particle approx.}}.
    \end{equation}
    Term~(I) involves two exact (continuous) posteriors, so the TV bound of Lemma~\ref{lem:Q_lipschitz_tv} gives $d_Q\le 2V_{\max}\|b_t^{\theta^\star}-b_t^{\bar\theta_{1:t}}\|_{\mathrm{TV}}$, which is bounded by the filter mismatch analysis below.

    \paragraph{Term (II): test-function bound via Feynman-Kac.}
    Conditional on $(a_{1:T},o_{1:T})$, the particle filter is a Feynman-Kac system with kernels $M_s(\cdot\mid s'):=P_{\bar\theta_s}(\cdot\mid s',a_s)$ and potentials $G_s(s'):=Z_{\bar\theta_s}(o_s\mid s')$. Under Assumption~\ref{aspt:particle_error}, the kernels satisfy $M_s(s'\mid x)\ge\varepsilon_{\mathrm{mix}}\,\nu(s')$ and the potentials satisfy $g_{\min}\le G_s(x)\le g_{\max}$ uniformly. We invoke the following time-uniform particle-filter bound. 
    
    \begin{lem}[{\cite[Lemma~2.3]{karjalainen23}; see also \cite[Theorem~7.4.4]{delMoral04}}]\label{lem:karjalainen}
    Let $(\eta_n)_{n\ge 0}$ and $(\eta_n^N)_{n\ge 0}$ be the exact and $N$-particle Feynman-Kac measures on a measurable space $E$. Suppose the Markov kernels $(M_\ell)_{\ell\ge 1}$ admit densities w.r.t.\ a common measure $\lambda$ satisfying $\underline M\le M_\ell(x,y)\le\overline M$ for all $\ell,x,y$, and the potentials satisfy $\underline G\le G_\ell(x)\le\overline G$ for all $\ell,x$. Then for all $n\ge 0$, $N\ge 2$, $p\ge 1$, and $\varphi$ with $\mathrm{osc}(\varphi)\le 1$:
    \begin{equation}\label{eq:karjalainen_lp}
    \bigl(\E\bigl[|\eta_n^N(\varphi)-\eta_n(\varphi)|^p\bigr]\bigr)^{1/p} 
    \;\le\;
    \frac{2\,d(p)^{1/p}}{\sqrt{N}}\,\frac{(\overline M/\underline M)^3\,\overline G}{\underline G},
    \end{equation}
    where $d(p)^{1/p}=O(\sqrt{p})$. The bound is uniform over $n$.
    \end{lem}

    Our Assumption~\ref{aspt:particle_error} implies the conditions of Lemma~\ref{lem:karjalainen} with $\underline M=\varepsilon_{\mathrm{mix}}$, $\overline M=1$, $\underline G=g_{\min}$, $\overline G=g_{\max}$.

    \begin{prop}[Particle filter $L^p$ bound under Assumption~\ref{aspt:particle_error}]\label{cor:pf_test_fn}
    Under Assumption~\ref{aspt:particle_error}, for every measurable $f:\gS\to\R$ with $\|f\|_\infty\le M$, all $p\ge 1$, and all $t\ge 0$:
    \begin{equation}\label{eq:delmoral_lp}
    \bigl(\E\bigl[|\eta_t^J(f)-\eta_t(f)|^p\bigr]\bigr)^{1/p}
    \;\le\;
    \frac{4\,d(p)^{1/p}\,M\,g_{\max}}{\varepsilon_{\mathrm{mix}}^3\,g_{\min}\,\sqrt{J}},
    \end{equation}
    where $d(p)^{1/p}=O(\sqrt{p})$, so $\eta_t^J(f)-\eta_t(f)$ is sub-Gaussian with parameter $\sigma^2=O(M^2 g_{\max}^2/(\varepsilon_{\mathrm{mix}}^6 g_{\min}^2 J))$. The bound is uniform over $t$.
    \end{prop}

    \begin{proof}
    By Lemma~\ref{lem:karjalainen} with $\underline M=\varepsilon_{\mathrm{mix}}$, $\overline M=1$, $\underline G=g_{\min}$, $\overline G=g_{\max}$, for any $\varphi$ with $\mathrm{osc}(\varphi)\le 1$: 
    \[
    \bigl(\E\bigl[|\eta_t^J(\varphi)-\eta_t(\varphi)|^p\bigr]\bigr)^{1/p}
    \;\le\;
    \frac{2\,d(p)^{1/p}\,g_{\max}}{\varepsilon_{\mathrm{mix}}^3\,g_{\min}\,\sqrt{J}},
    \]
    uniformly over $t$. For $f$ with $\|f\|_\infty\le M$, applying this to $\varphi:=f/(2M)$ (which has $\mathrm{osc}(\varphi)\le 1$) and rescaling by $2M$ yields~\eqref{eq:delmoral_lp}. Sub-Gaussianity follows from the standard moment-to-tail conversion.
    \end{proof}

    \begin{lem}[Unconditional bound under adaptive actions]\label{lem:adaptive_conditioning}
    Let $(a_s,o_s)_{s\ge 1}$ be an adapted action-observation process where $a_s=a_s(b_s^J)$ may depend on the particle system. Under Assumption~\ref{aspt:particle_error}, for every measurable $f:\gS\to\R$ with $\|f\|_\infty\le M$, all $p\ge 1$, and all $t\ge 0$:
    \begin{equation}\label{eq:unconditional_lp}
    \bigl(\E\bigl[|b_t^{J,\bar\theta_{1:t}}(f)-b_t^{\bar\theta_{1:t}}(f)|^p\bigr]\bigr)^{1/p}
    \;\le\;
    \frac{4\,d(p)^{1/p}\,M\,g_{\max}}{\varepsilon_{\mathrm{mix}}^3\,g_{\min}\,\sqrt{J}}.
    \end{equation}
    \end{lem}

    \begin{proof}
    Let $\gH_T:=(a_{1:T},o_{1:T},\bar\theta_{1:T})$. Conditional on $\gH_T$, the Feynman-Kac kernels and potentials are deterministic, so Proposition~\ref{cor:pf_test_fn} gives
    \[
    \E\bigl[|b_t^J(f)-b_t(f)|^p\,\bigm|\,\gH_T\bigr]
    \;\le\;
    \left(\frac{4\,d(p)^{1/p}\,M\,g_{\max}}{\varepsilon_{\mathrm{mix}}^3\,g_{\min}\,\sqrt{J}}\right)^p
    =:C_p^p.
    \]
    Since $C_p$ is independent of $\gH_T$, the tower property yields $\E[|b_t^J(f)-b_t(f)|^p]\le C_p^p$.
    \end{proof}

    By Lemma~\ref{lem:Q_lipschitz} with truncation depth $H_{\mathrm{pf}}$:
    \begin{equation}\label{eq:dQ_via_test}
    d_Q(b_t^{\bar\theta_{1:t}},b_t^{J,\bar\theta_{1:t}})
    \;\le\;
    \sup_{\theta,a}\sum_{\ell=0}^{H_{\mathrm{pf}}-1}\gamma^\ell C_{\mathrm{BU}}^\ell\sup_{f\in\gF_\ell(\theta,a)}|b_t(f)-b_t^J(f)|
    +2\gamma^{H_{\mathrm{pf}}} V_{\max}.
    \end{equation}
    We optimize $H_{\mathrm{pf}}$ to balance the sum and tail below.

    \emph{Pointwise bound.} For $f\in\gF_\ell(\theta,a)$ with $\|f\|_\infty\le g_{\max}^\ell V_{\max}$, Lemma~\ref{lem:adaptive_conditioning} gives
    \begin{equation}\label{eq:pointwise_pf}
    \bigl(\E\bigl[|b_t^J(f)-b_t(f)|^p\bigr]\bigr)^{1/p}
    \;\le\;
    \frac{4\,d(p)^{1/p}\,g_{\max}^{\ell+1}\,V_{\max}}{\varepsilon_{\mathrm{mix}}^3\,g_{\min}\,\sqrt{J}}
    =:\sigma_\ell\,d(p)^{1/p}.
    \end{equation}
    Since $d(p)^{1/p}\le C_0\sqrt{p}$ (Corollary~\ref{cor:pf_test_fn}), $b_t^J(f)-b_t(f)$ is sub-Gaussian with parameter $\sigma_\ell^2$.

    \emph{Covering-number argument.} For $N$ sub-Gaussian random variables with common parameter $\sigma^2$:
    \begin{equation}\label{eq:max_ineq}
    \E\bigl[\max_{i\le N}|X_i|\bigr]
    \;\le\;
    \sigma\sqrt{2\log(2N)}.
    \end{equation} 
    At depth $\ell$, $\bigcup_{\theta,a}\gF_\ell(\theta,a)$ is covered by $N_\ell:=\gN_\Theta(1/T)\cdot|\gA|^\ell$ functions (after discretizing $\Theta$ at resolution $1/T$, approximation error $O(1/T)$). Applying~\eqref{eq:max_ineq} and summing over depths with weights $\gamma^\ell C_{\mathrm{BU}}^\ell$:
    \begin{align}
    \E\bigl[d_Q(b_t^{\bar\theta_{1:t}},b_t^{J,\bar\theta_{1:t}})\bigr]
    &\le\sum_{\ell=0}^{H_{\mathrm{pf}}-1}\gamma^\ell C_{\mathrm{BU}}^\ell\,\sigma_\ell\sqrt{2\log(2N_\ell)}+2\gamma^{H_{\mathrm{pf}}} V_{\max}+O(1/T)
    % &=\frac{1}{\sqrt{J}}\sum_{\ell=0}^{H_{\mathrm{pf}}-1}\underbrace{\gamma^\ell\left(\frac{2g_{\max}}{g_{\min}^2}\right)^{\ell}}_{\text{amplification}}\cdot\underbrace{\frac{4\,g_{\max}^{\ell+1}\,V_{\max}}{\varepsilon_{\mathrm{mix}}^3\,g_{\min}}}_{\sigma_\ell\sqrt{J}}\cdot\sqrt{2\log(2N_\ell)}+O(1/\sqrt{J}). 
    \label{eq:Cpf_prime_def}
    \end{align} 
    Define $\alpha:=2\gamma g_{\max}^2/g_{\min}^2$, $r:=\alpha/\gamma=2g_{\max}^2/g_{\min}^2$, and $c_1:=4g_{\max}V_{\max}/(\varepsilon_{\mathrm{mix}}^3 g_{\min})$, so that $\sigma_\ell=c_1\,g_{\max}^\ell/\sqrt{J}$ and the $\ell$-th term equals 
\[ 
\frac{c_1}{\sqrt{J}}\,\alpha^\ell\,\sqrt{2\log\bigl(2\gN_\Theta(1/T)\,|\gA|^\ell\bigr)}. 
\]
Since $\sqrt{\log(ab)}\le\sqrt{\log a}+\sqrt{\log b}$, the sum in~\eqref{eq:Cpf_prime_def} is bounded by
\begin{equation}\label{eq:pf_balance}
\E\bigl[d_Q^{\mathrm{(II)}}\bigr]
\;\le\;
\frac{c_1\sqrt{2}}{\sqrt{J}}\,\Biggl( 
\sqrt{\log(2\gN_\Theta(1/T))}\sum_{\ell=0}^{H_{\mathrm{pf}}-1}\alpha^\ell
+\sqrt{\log|\gA|}\sum_{\ell=0}^{H_{\mathrm{pf}}-1}\alpha^\ell\sqrt{\ell}
\Biggr)
+2\gamma^{H_{\mathrm{pf}}} V_{\max}.
\end{equation}
We bound the two sums in two cases.

\emph{Case 1: $\alpha\ge 1$.} Choose
\[
H_{\mathrm{pf}}=\left\lceil\frac{\log(4V_{\max}\sqrt{J})}{\log r}\right\rceil
=\left\lceil\frac{\log(4V_{\max}\sqrt{J})}{\log(2g_{\max}^2/g_{\min}^2)}\right\rceil.
\]
Using $\sum_{\ell=0}^{H-1}\alpha^\ell\le H\alpha^{H-1}$ and $\sum_{\ell=0}^{H-1}\alpha^\ell\sqrt{\ell}\le H^{3/2}\alpha^{H-1}$ (since $\alpha\ge 1$):
\begin{equation}\label{eq:sum_alpha_ge1}
\text{sum}\;\le\;\frac{c_1\sqrt{2}\,H_{\mathrm{pf}}^{3/2}\,\alpha^{H_{\mathrm{pf}}-1}}{\sqrt{J}}\cdot\sqrt{\log\bigl(2\gN_\Theta(1/T)\,|\gA|^{H_{\mathrm{pf}}}\bigr)}.
\end{equation}
Writing $\alpha^{H_{\mathrm{pf}}-1}=\alpha^{H_{\mathrm{pf}}}/\alpha=\gamma^{H_{\mathrm{pf}}}r^{H_{\mathrm{pf}}}/\alpha$ and substituting $r^{H_{\mathrm{pf}}}=4V_{\max}\sqrt{J}$:
\begin{equation}\label{eq:sum_case1_final}
\text{sum}\;\le\;\frac{4\sqrt{2}\,c_1\,V_{\max}\,H_{\mathrm{pf}}^{3/2}}{\alpha}\cdot\sqrt{\log\bigl(2\gN_\Theta(1/T)\,|\gA|^{H_{\mathrm{pf}}}\bigr)}\cdot\gamma^{H_{\mathrm{pf}}}.
\end{equation}
For the tail:
\begin{equation}\label{eq:tail_case1}
\gamma^{H_{\mathrm{pf}}}=\gamma^{\log(4V_{\max}\sqrt{J})/\log r}=(4V_{\max}\sqrt{J})^{-\log(1/\gamma)/\log r}=(4V_{\max})^{-1/\kappa_{\mathrm{pf}}}\cdot J^{-\frac{1}{2\kappa_{\mathrm{pf}}}},
\end{equation}
where $\kappa_{\mathrm{pf}}=\log r/\log(1/\gamma)>1$. Substituting $c_1=4g_{\max}V_{\max}/(\varepsilon_{\mathrm{mix}}^3 g_{\min})$, $H_{\mathrm{pf}}=\lceil\log(4V_{\max}\sqrt{J})/\log r\rceil$, and using $\log(2\gN_\Theta(1/T)\,|\gA|^{H_{\mathrm{pf}}})\le\log(2\gN_\Theta(1/T))+H_{\mathrm{pf}}\log|\gA|$, the total (sum $+$ tail) is
\begin{align}
\E\bigl[d_Q^{\mathrm{(II)}}\bigr]
&\;\le\;
\Biggl(\frac{16\sqrt{2}\,g_{\max}\,V_{\max}^2}{\alpha\,\varepsilon_{\mathrm{mix}}^3\,g_{\min}}\cdot \left(\frac{\log(4V_{\max}\sqrt{J})}{\log r}\right)^{3/2}\sqrt{\log\bigl(2\gN_\Theta(1/T)\bigr)+\frac{\log(4V_{\max}\sqrt{J})}{\log r}\log|\gA|}\notag\\
&\qquad\qquad+2V_{\max}\Biggr)\cdot(4V_{\max})^{-1/\kappa_{\mathrm{pf}}}\cdot J^{-\frac{1}{2\kappa_{\mathrm{pf}}}}.\label{eq:case1_full}
\end{align}

\emph{Case 2: $\alpha<1$.} Choose
\[
H_{\mathrm{pf}}=\left\lceil\frac{\log(4V_{\max}\sqrt{J})}{\log(1/\gamma)}\right\rceil.
\]
The geometric series converges: $\sum_{\ell\ge 0}\alpha^\ell=1/(1-\alpha)$ and $\sum_{\ell\ge 0}\alpha^\ell\sqrt{\ell}\le c_2/(1-\alpha)^{3/2}$ for a universal $c_2>0$. Hence
\[
\text{sum}\;\le\;\frac{c_1\,c_2\,\sqrt{2}}{(1-\alpha)^{3/2}\,\sqrt{J}}\cdot\sqrt{\log\bigl(2\gN_\Theta(1/T)\bigr)+\log|\gA|}
\;=\;O\!\left(J^{-1/2}\right).
\]
For the tail: $\gamma^{H_{\mathrm{pf}}}=(4V_{\max}\sqrt{J})^{-1}=O(J^{-1/2})$. Here $\kappa_{\mathrm{pf}}=1$, and both terms are $O(J^{-1/2})$.

\emph{Unified bound.} The choice~\eqref{eq:Hpf_choice}, $H_{\mathrm{pf}}=\lceil\log(4V_{\max}\sqrt{J})/\log\max(1/\gamma,r)\rceil$, covers both cases and yields
\[
\E\bigl[d_Q^{\mathrm{(II)}}\bigr]\;\le\;C_{\mathrm{pf}}\cdot J^{-\frac{1}{2\kappa_{\mathrm{pf}}}},
\qquad
\kappa_{\mathrm{pf}}=\max\!\left(1,\frac{\log r}{\log(1/\gamma)}\right),
\]
where
\begin{equation}\label{eq:Cpf_explicit}
C_{\mathrm{pf}}:=\frac{c_3\,V_{\max}\,g_{\max}\,H_{\mathrm{pf}}^{3/2}}{\varepsilon_{\mathrm{mix}}^3\,g_{\min}}\cdot\sqrt{\log\bigl(2\gN_\Theta(1/T)\,|\gA|^{H_{\mathrm{pf}}}\bigr)},
\end{equation}
with $c_3$ depending on $\alpha$ and $g_{\max}/g_{\min}$. In particular, the $|\gA|$ dependence enters through $\sqrt{H_{\mathrm{pf}}\log|\gA|}$.

\paragraph{Term (I): filter mismatch.}
By Lemma~\ref{lem:Q_lipschitz_tv},
\[
d_Q(b_t^{\theta^\star},b_t^{\bar\theta_{1:t}})\;\le\; 2V_{\max}\|b_t^{\theta^\star}-b_t^{\bar\theta_{1:t}}\|_{\mathrm{TV}}.
\]
Define $\delta_s:=\sup_{s'}\|P_{\bar\theta_s}(\cdot\mid s',a_s)-P_{\theta^\star}(\cdot\mid s',a_s)\|_{\mathrm{TV}}+\sup_{s'}|Z_{\bar\theta_s}(o_s\mid s')/Z_{\theta^\star}(o_s\mid s')-1|$ and the predict-update kernel
\[
K_\theta^{a,o}(s,ds'):=\frac{Z_\theta(o\mid s')P_\theta(s'\mid s,a)}{\int Z_\theta(o\mid s'')P_\theta(s''\mid s,a)\,ds''}\,ds'.
\]
Under Assumption~\ref{aspt:particle_error}, bounding the numerator below and denominator above:
\begin{equation}\label{eq:kernel_doeblin}
K_\theta^{a,o}(s,ds')\;\ge\;\frac{g_{\min}\varepsilon_{\mathrm{mix}}}{g_{\max}}\,\nu(s')\,ds'
\qquad\forall\,s\in\gS,
\end{equation}
giving Dobrushin coefficient $\rho:=1-\varepsilon_{\mathrm{mix}}g_{\min}/g_{\max}$. Standard Dobrushin contraction~\citep[Theorem~4.3.1]{doucmoulinesolsson2014} under different models $\theta^\star,\bar\theta_s$ yields
\begin{equation}\label{eq:filter_recursion}
\|b_s^{\theta^\star}-b_s^{\bar\theta_{1:s}}\|_{\mathrm{TV}}
\;\le\;
\rho\,\|b_{s-1}^{\theta^\star}-b_{s-1}^{\bar\theta_{1:s-1}}\|_{\mathrm{TV}}+C_{\mathrm{filt}}\,\delta_s,
\qquad C_{\mathrm{filt}}:=g_{\max}/g_{\min}.
\end{equation}
Unrolling and using $\delta_s\le d_{\mathrm{model}}^{a_s}(\bar\theta_s,\theta^\star)$:
\[
\|b_t^{\theta^\star}-b_t^{\bar\theta_{1:t}}\|_{\mathrm{TV}}
\;\le\;
C_{\mathrm{filt}}\sum_{s=1}^t \rho^{t-s}\,d_{\mathrm{model}}^{a_s}(\bar\theta_s,\theta^\star).
\]

\paragraph{Combining.}
Substituting into~\eqref{eq:pf_triangle}:
\[
\E\,d_Q(b_t,b_t^J)
\;\le\;
\frac{C_{\mathrm{pf}}'}{\sqrt{J}}+2V_{\max}\,C_{\mathrm{filt}}\sum_{s=1}^t\rho^{t-s}\,d_{\mathrm{model}}^{a_s}(\bar\theta_s,\theta^\star).
\]
When summed over $t$, exchanging summation order gives
$\sum_{t=1}^T\sum_{s=1}^t\rho^{t-s}\,d_{\mathrm{model}}^{a_s}
\le\frac{1}{1-\rho}\sum_{s=1}^T d_{\mathrm{model}}^{a_s}.$
\end{proof}


%% ===================================================================
\subsection{Model Estimation Error (Terms A + C) and Filter Mismatch}
\label{app:model_error_proofs}

\paragraph{Truncation.}
Define $V_\theta^{\pi,B,H}(b):=\E_\theta^\pi[\sum_{u=0}^{H-1}\gamma^u r_u\mid b_0=b]$. Since $|V_\theta^{\pi,B}-V_\theta^{\pi,B,H}|\le \gamma^H V_{\max}$:
\begin{equation}\label{eq:truncation_error}
\text{(A)}_t+\text{(C)}_t
\;\le\;
|V_{\theta^\star}^{\star,B,H}(b_t)-V_{\bar\theta_{k(t)}}^{\star,B,H}(b_t)|
+|V_{\bar\theta_{k(t)}}^{\pi_t,B,H}(b_t)-V_{\theta^\star}^{\pi_t,B,H}(b_t)|
+\frac{4\gamma^H R_{\max}}{1-\gamma}.
\end{equation}
Choose $H=\lceil\log(2R_{\max}T/(1-\gamma))/\log(1/\gamma)\rceil$ so the tail contributes $O(1)$.

\paragraph{Confidence sets.}
Let $\PP_{\theta,H}^\pi$ denote the length-$H$ trajectory distribution under $(\theta,\pi)$. Define
\begin{equation}\label{eq:conf_set}
\gC_k:=\bigl\{\theta\in\Theta:\gL_{T_k}(\theta)\ge \gL_{T_k}(\bar\theta_k)-\beta\bigr\},
\end{equation}
where $\gL_{T_k}$ is the log-likelihood over epochs $1,\ldots,k-1$. By construction $\bar\theta_k\in\gC_k$, and $\theta^\star\in\gC_k$ for all $k$ w.h.p.\ when $\beta\ge c_0\log[\gN_\Theta(1/T)T/\delta_{\mathrm{opt}}]$~\citep[Prop.~B.1]{liu2023optimistic}.

\begin{prop}[Likelihood-to-$\ell_1$ concentration]\label{prop:lik_to_l1}
There exists $c>0$ such that w.p.\ $\ge 1-\delta$, simultaneously for all $n\in[T]$ and all $\theta\in\Theta$ with $\gL_n(\theta)\ge\gL_n(\theta^\star)-\beta$:
\begin{equation}\label{eq:lik_to_l1}
\sum_{i=1}^{n}\bigl\|\PP_{\theta,H}^{\pi^i}-\PP_{\theta^\star,H}^{\pi^i}\bigr\|_1^2
\;\le\;
c\,\beta_T,
\qquad
\beta_T:=\beta+\log[\gN_\Theta(1/T)\,T/\delta].
\end{equation}
\end{prop}

\begin{proof}
The episodic version is given by~\citet[Prop.~B.2]{liu2023optimistic}. We transfer to the online setting via a blocking argument: partition the trajectory into $\lfloor T/(2H)\rfloor$ blocks of length $2H$, each split into a burn-in half and a data half. Choose $H$ so that
\[
\rho^H = (1-\varepsilon_{\mathrm{mix}}g_{\min}/g_{\max})^H \;\le\; 1/T^2,
\qquad\text{i.e.,}\quad
H \;\ge\; \frac{2\log T}{\log(1/\rho)}.
\]
Our choice $H=O(\log T/\log(1/\gamma))$ satisfies this when $\varepsilon_{\mathrm{mix}}g_{\min}/g_{\max}=\Omega(1-\gamma)$; otherwise, increase $H$ by a constant factor. With $\rho^H\le 1/T^2$, the data halves are approximately independent (joint TV perturbation $O(1/T)$). Applying the episodic bound to the data halves gives~\eqref{eq:lik_to_l1} up to a constant factor.
\end{proof}

\begin{cor}[Epoch-level MLE concentration]\label{cor:epoch_mle}
Under the same conditions as Proposition~\ref{prop:lik_to_l1}, w.p.\ $\ge 1-\delta$, simultaneously for all epochs $k\in[K]$:
\begin{equation}\label{eq:hist_l1}
\sum_{i=1}^{T_{k+1}-1}\bigl\|\PP_{\bar\theta_k,H}^{\pi^i}-\PP_{\theta^\star,H}^{\pi^i}\bigr\|_1^2
\;\le\;c\,\beta_T,
\end{equation}
where $c>0$ is a universal constant (at most twice the constant in Proposition~\ref{prop:lik_to_l1}).
\end{cor}

\begin{proof}
Fix $k\in[K]$ and write $\|\PP_\theta^\pi\|_1:=\|\PP_{\theta,H}^\pi-\PP_{\theta^\star,H}^\pi\|_1$.

\emph{Historical data ($i\le T_k$).}
By MLE optimality, $\bar\theta_k$ lies in $\gC_k$, so Proposition~\ref{prop:lik_to_l1} at $n=T_k$ gives
\begin{equation}\label{eq:hist_bound_epoch_start}
\sum_{i=1}^{T_k}\bigl\|\PP_{\bar\theta_k}^{\pi^i}\bigr\|_1^2\;\le\;c\,\beta_T.
\end{equation}

\emph{Within-epoch data ($i\in E_k$).}
Since $\bar\theta_k$ is $\gF_{T_k}$-measurable, we use a self-normalized argument for the fixed hypothesis $\bar\theta_k$. Define
\[
L_n^{(k)}:=\sum_{i=T_k+1}^{T_k+n}\bigl[\ell_i(\theta^\star)-\ell_i(\bar\theta_k)\bigr],
\qquad
M_n^{(k)}:=\sum_{i=T_k+1}^{T_k+n}\bigl[\mathrm{KL}(\PP_{\theta^\star}^{\pi^i}\|\PP_{\bar\theta_k}^{\pi^i})-\bigl(\ell_i(\theta^\star)-\ell_i(\bar\theta_k)\bigr)\bigr],
\]
where $\ell_i(\theta):=\log p_\theta(o_i\mid s_i,a_i)$. By Pinsker's inequality and blocking:
\begin{equation}\label{eq:pinsker_within}
\frac{1}{2}\sum_{i\in E_k}\bigl\|\PP_{\bar\theta_k}^{\pi^i}\bigr\|_1^2
\;\le\;
L_{|E_k|}^{(k)}+M_{|E_k|}^{(k)}.
\end{equation}
$M_n^{(k)}$ is a supermartingale (each increment has non-positive conditional expectation). The maximal inequality and a covering-number union bound give
\[
\Pr\!\left(\sup_{n\le|E_k|}M_n^{(k)}\ge c_1\log\bigl[\gN_\Theta(1/T)/\delta\bigr]\right)\le\delta/K.
\]
For $L_{|E_k|}^{(k)}$: each $|\ell_i(\theta^\star)-\ell_i(\bar\theta_k)|\le\log(g_{\max}/g_{\min})$ and $\E[\ell_i(\theta^\star)-\ell_i(\bar\theta_k)\mid\gF_{i-1}]\ge 0$. By Azuma-Hoeffding:
\[
\Pr\!\left(L_{|E_k|}^{(k)}\ge c_2\sqrt{|E_k|\log(1/\delta)}\log(g_{\max}/g_{\min})\right)\le\delta/K.
\]
Since $\sqrt{T}\log(g_{\max}/g_{\min})$ is absorbed into $\beta_T$:
\begin{equation}\label{eq:within_epoch_bound}
\sum_{i\in E_k}\bigl\|\PP_{\bar\theta_k}^{\pi^i}\bigr\|_1^2
\;\le\;c\,\beta_T.
\end{equation}

\emph{Combining.} Summing~\eqref{eq:hist_bound_epoch_start} and~\eqref{eq:within_epoch_bound}, redefining $c\leftarrow 2c$, and taking a union bound over $k\in[K]$ gives~\eqref{eq:hist_l1}.
\end{proof}

Transferring model error from historical to arbitrary policies requires action-space coverage via $\varepsilon$-greedy exploration.

\begin{lem}[Trajectory coverage under Doeblin mixing and $\varepsilon$-greedy]\label{lem:coverage}
Under Assumption~\ref{aspt:particle_error}, for any two models $\theta,\theta'\in\Theta$, any policy $\pi$, and any policy $\pi_\varepsilon$ satisfying $\pi_\varepsilon(a\mid h)\ge\varepsilon/|\gA|$ for all $a,h$ (with $\varepsilon\in(0,1]$),
\begin{equation}\label{eq:coverage_bound}
\bigl\|\PP_{\theta,H}^\pi-\PP_{\theta',H}^\pi\bigr\|_1
\;\le\;
C_{\mathrm{cov}}\,
\bigl\|\PP_{\theta,H}^{\pi_\varepsilon}-\PP_{\theta',H}^{\pi_\varepsilon}\bigr\|_1,
\qquad
C_{\mathrm{cov}}:=\frac{H\,|\gA|}{\varepsilon}\cdot\frac{g_{\max}^2}{\varepsilon_{\mathrm{mix}}\,g_{\min}^2}.
\end{equation}
\end{lem}

\begin{proof}
\emph{Step 1: Telescoping.}
Let $\PP^{(k)}$ be the hybrid using $\theta$ for steps $1,\ldots,k$ and $\theta'$ for steps $k+1,\ldots,H$, under policy $\pi$. Let $\Delta_k(s,a):=\|P_\theta(\cdot\mid s,a)-P_{\theta'}(\cdot\mid s,a)\|_{\mathrm{TV}}$ and $\Delta^Z(s):=\sup_o|Z_\theta(o\mid s)/Z_{\theta'}(o\mid s)-1|$. Then
\begin{equation}\label{eq:telescope}
\|\PP_{\theta,H}^\pi-\PP_{\theta',H}^\pi\|_1
\;\le\;
\sum_{k=0}^{H-1}\|\PP^{(k+1)}-\PP^{(k)}\|_1,
\end{equation}
\begin{equation}\label{eq:per_step_bound}
\|\PP^{(k+1)}-\PP^{(k)}\|_1
\;\le\;
\E_{\mu_k^\pi}\!\bigl[\sup_{a}\Delta_k(S_k,a)+\Delta^Z(S_k)\bigr],
\end{equation}
where $\mu_k^\pi$ is the state marginal at step $k$ under the hybrid.

\emph{Step 2: State-marginal ratio.}
By the Doeblin decomposition $\tilde\mu_k^\pi \ge \varepsilon_{\mathrm{mix}}\nu$ and $\tilde\mu_k^{\pi'}\ge\varepsilon_{\mathrm{mix}}\nu$, for any $h:\gS\to[0,1]$:
\[
\E_{\tilde\mu_k^\pi}[h]
\;\le\; \varepsilon_{\mathrm{mix}}\E_\nu[h] + (1-\varepsilon_{\mathrm{mix}})
\;\le\; \frac{1}{\varepsilon_{\mathrm{mix}}}\E_{\tilde\mu_k^{\pi'}}[h].
\]
Including the Bayes update normalization ($g_{\max}/g_{\min}$):
\[
\E_{\mu_k^\pi}[h]\;\le\; \frac{g_{\max}}{g_{\min}\varepsilon_{\mathrm{mix}}}\,\E_{\mu_k^{\pi'}}[h].
\]

\emph{Step 3: Action coverage.}
Since $\pi_\varepsilon(a\mid h)\ge\varepsilon/|\gA|$:
\[
\sup_{a}\Delta_k(s,a)
\;\le\;
\frac{|\gA|}{\varepsilon}\,\E_{A\sim\pi_\varepsilon}[\Delta_k(s,A)].
\]

\emph{Step 4: Assembly.}
Combining Steps 1--3 term-by-term:
\[
\|\PP^{(k+1)}-\PP^{(k)}\|_1
\;\le\;
\frac{|\gA|}{\varepsilon}\cdot\frac{g_{\max}^2}{\varepsilon_{\mathrm{mix}}\,g_{\min}^2}\cdot\|\PP^{(k+1)}_{\pi_\varepsilon}-\PP^{(k)}_{\pi_\varepsilon}\|_1.
\]
Summing over $k=0,\ldots,H-1$ yields~\eqref{eq:coverage_bound}.\qedhere
\end{proof}
 
% \begin{rmk}
% Without the Doeblin condition, the coverage constant would scale as $(|\gA|/\varepsilon)^H$ from a pure propensity-score argument (bounding $\prod_t\pi(a_t)/\pi_\varepsilon(a_t)$). The Doeblin condition reduces this to $O(H\,|\gA|/\varepsilon)$ because the state-space mixing ensures that the effect of each action choice does not compound multiplicatively across steps. Since $H=O(\log T)$, the coverage constant is $O(|\gA|\log T/\varepsilon)$ --- logarithmic in $T$ and linear in $|\gA|/\varepsilon$.
% \end{rmk}

\begin{lem}[Simulation lemma]\label{lem:simulation}
For any $\theta,\theta'\in\Theta$, $\pi$, and $b$:
\[
\bigl|V_\theta^{\pi,B,H}(b)-V_{\theta'}^{\pi,B,H}(b)\bigr| 
\;\le\;
\frac{R_{\max}}{1-\gamma}\,\bigl\|\PP_{\theta,H}^\pi-\PP_{\theta',H}^\pi\bigr\|_1.
\]
\end{lem}
\begin{proof}
Telescope the $H$-truncated value difference into one-step terms, each bounded by $R_{\max}$.
\end{proof}

\begin{prop}[Control of (A) + (C) + filter mismatch]\label{prop:model_error_combined}
Under Assumption~\ref{aspt:particle_error}, consider the doubling-epoch variant of Algorithm~\ref{alg:pf_mcts} in which time is partitioned into epochs $E_k=[T_k,T_{k+1})$ with $T_k=2^{k-1}$, $k=1,\ldots,K:=\lceil\log_2 T\rceil$. Within each epoch $E_k$, the model $\bar\theta_k$ is frozen (computed from data in epochs $1,\ldots,k-1$) and the agent uses $\varepsilon_k$-greedy exploration with a constant rate $\varepsilon_k$. Let $K_{\mathrm{cov}}:=H|\gA|\,g_{\max}^2/(\varepsilon_{\mathrm{mix}}\,g_{\min}^2)$. Then with probability $\ge 1-\delta_{\mathrm{opt}}$,
\begin{align}
&\sum_{t=1}^T\text{(C)}_t
\;\le\;
\frac{R_{\max}}{1-\gamma}\sqrt{c\,\beta_T}\sum_{k=1}^K\sqrt{|E_k|}+\frac{2\gamma^H R_{\max} T}{1-\gamma}, \label{eq:c_bound}\\
&\sum_{t=1}^T\text{(A)}_t
\;\le\;
\frac{R_{\max}\,K_{\mathrm{cov}}}{1-\gamma}\sqrt{c\,\beta_T}\sum_{k=2}^K\frac{|E_k|}{\sqrt{\sum_{j=1}^{k-1}\varepsilon_j^2|E_j|}}+\frac{2\gamma^H R_{\max} T}{1-\gamma}, \label{eq:a_bound}\\
&\sum_{t=1}^T d_{\mathrm{model}}^{a_t}(\bar\theta_{k(t)},\theta^\star)
\;\le\;
C_{\mathrm{mix}}\sqrt{c\,\beta_T}\sum_{k=1}^K\sqrt{|E_k|}, \label{eq:dmodel_bound}
\end{align}
where $d_{\mathrm{model}}^{a}(\theta,\theta'):=\sup_{s\in\gS}\|P_\theta(\cdot\mid s,a)-P_{\theta'}(\cdot\mid s,a)\|_{\mathrm{TV}}+\sup_{s\in\gS,o\in\gO}\bigl|Z_\theta(o\mid s)/Z_{\theta'}(o\mid s)-1\bigr|$ is the per-action model discrepancy, and $C_{\mathrm{mix}}$ depends on $\varepsilon_{\mathrm{mix}},g_{\min},g_{\max},H$.
\end{prop}



\begin{proof}
By Corollary~\ref{cor:epoch_mle}, equation~\eqref{eq:hist_l1} holds for all $k\le K$.
For brevity, write $\|\PP_{\theta}^{\pi}\|_1:=\|\PP_{\theta,H}^{\pi}-\PP_{\theta^\star,H}^{\pi}\|_1$ throughout this proof.

\emph{Term (C).} By Lemma~\ref{lem:simulation},
\[
\text{(C)}_t\;\le\; \frac{R_{\max}}{1-\gamma}\bigl\|\PP_{\bar\theta_{k(t)}}^{\pi^t}\bigr\|_1+\frac{2\gamma^H R_{\max}}{1-\gamma}. 
\]
Since $\bar\theta_k$ is fixed within $E_k$, Cauchy-Schwarz and~\eqref{eq:hist_l1} give 
\[
\sum_{t\in E_k}\bigl\|\PP_{\bar\theta_k}^{\pi^t}\bigr\|_1
\;\le\;\sqrt{|E_k|\sum_{t\in E_k}\bigl\|\PP_{\bar\theta_k}^{\pi^t}\bigr\|_1^2}
\;\le\;\sqrt{|E_k|\cdot c\,\beta_T}.
\]
Summing across epochs: $\sum_{k=1}^K\sqrt{|E_k|\cdot c\beta_T}=\sqrt{c\beta_T}\sum_k 2^{(k-1)/2}=O(\sqrt{T\beta_T})$.

\emph{Term (A).} By Lemmas~\ref{lem:simulation} and~\ref{lem:coverage}:
\[
\text{(A)}_t\;\le\;\frac{R_{\max}}{1-\gamma}\sup_\pi\bigl\|\PP_{\bar\theta_{k(t)}}^{\pi}\bigr\|_1+\frac{2\gamma^H R_{\max}}{1-\gamma},
\qquad
\bigl\|\PP_{\bar\theta_k}^{\pi}\bigr\|_1\;\le\;\frac{K_{\mathrm{cov}}}{\varepsilon_j}\bigl\|\PP_{\bar\theta_k}^{\pi^j}\bigr\|_1\quad\forall\,j<k.
\]
Squaring, multiplying by $\varepsilon_j^2$, and summing over prior epochs using~\eqref{eq:hist_l1}:
\[
\bigl\|\PP_{\bar\theta_k}^{\pi}\bigr\|_1^2\sum_{j<k}\varepsilon_j^2|E_j|
\;\le\; K_{\mathrm{cov}}^2\sum_{j<k}\sum_{t\in E_j}\bigl\|\PP_{\bar\theta_k}^{\pi^t}\bigr\|_1^2
\;\le\; K_{\mathrm{cov}}^2\,c\,\beta_T
\quad\implies\quad
\sup_\pi\bigl\|\PP_{\bar\theta_k}^{\pi}\bigr\|_1\;\le\;\frac{K_{\mathrm{cov}}\sqrt{c\beta_T}}{\sqrt{\sum_{j<k}\varepsilon_j^2|E_j|}}.
\]
Multiplying by $|E_k|$ and summing gives~\eqref{eq:a_bound}.

\emph{Filter mismatch.}
By the telescoping in Lemma~\ref{lem:coverage} (Steps~1--2 only, without the $|\gA|/\varepsilon$ action-coverage factor since we evaluate at the executed action $a_t$):
\[
d_{\mathrm{model}}^{a_t}(\bar\theta_{k(t)},\theta^\star)
\;\le\;
C_{\mathrm{mix}}\,\bigl\|\PP_{\bar\theta_{k(t)}}^{\pi^t}\bigr\|_1,
\qquad C_{\mathrm{mix}}:=\frac{H\,g_{\max}^2}{\varepsilon_{\mathrm{mix}}\,g_{\min}^2}.
\]
Cauchy-Schwarz and~\eqref{eq:hist_l1} give $\sum_{t\in E_k}\|\PP_{\bar\theta_k}^{\pi^t}\|_1\le\sqrt{|E_k|\cdot c\,\beta_T}$. Summing across epochs yields~\eqref{eq:dmodel_bound}.
\end{proof}


%% ===================================================================
\subsection{Proof of the Main Regret Bound (Theorem~\ref{thm:final_regret_bound})}
\label{app:regret_proof}

\begin{proof}[Proof of Theorem~\ref{thm:final_regret_bound}]
Use doubling epochs $E_k=[T_k,T_{k+1})$, $T_k=2^{k-1}$, $K=\lceil\log_2 T\rceil$, with $\varepsilon_k$-greedy exploration within each epoch. Decompose:
\[
\mathrm{Reg}_T \;\le\; \underbrace{\frac{2R_{\max}}{1-\gamma}\sum_k|E_k|\varepsilon_k}_{\text{exploration}} + \E\sum_{t=1}^T\bigl[\text{(A)}_t+\text{(B)}_t+\text{(C)}_t\bigr].
\]

\emph{Term (C).} By~\eqref{eq:c_bound} (choosing $\delta_{\mathrm{opt}}=1/T$):
\[
\E\sum_{t}\text{(C)}_t\;\le\; \frac{R_{\max}\sqrt{c\beta_T}}{1-\gamma}\sum_{k=1}^K\sqrt{|E_k|}+O(1)
\;=\;\widetilde O\!\left(\frac{\sqrt{T}}{1-\gamma}\right).
\]

\emph{Term (A).} By~\eqref{eq:a_bound} with $\varepsilon_k=\eta/\sqrt{|E_k|}$ (so $\varepsilon_k^2|E_k|=\eta^2$):
\[
\E\sum_t\text{(A)}_t\;\le\;\frac{R_{\max}K_{\mathrm{cov}}\sqrt{c\beta_T}}{(1-\gamma)\eta}\sum_{k=2}^K\frac{|E_k|}{\sqrt{k-1}}+O(1).
\]
Since $\sum_{k=2}^K 2^{k-1}/\sqrt{k-1}=\Theta(T/\sqrt{\log T})$, this scales as $K_{\mathrm{cov}}\sqrt{\beta_T}\cdot T/(\eta\sqrt{\log T}(1-\gamma))$.

\emph{Term (B).} By Theorem~\ref{thm:planning_state_bound}, Propositions~\ref{prop:structural_error},~\ref{prop:mc_concentration},~\ref{prop:particle_error}:
\[
\E\sum_{t}\text{(B)}_t
\;\le\; \frac{2C_{\mathrm{pf}}\,T}{(1-\gamma)}\cdot J^{-\frac{1}{2\kappa_{\mathrm{pf}}}}
+\frac{C_{\mathrm{mc}}V_{\max}T}{(1-\gamma)\sqrt{N}}
+\frac{2\gamma^{d+1}}{(1-\gamma)^2}\sum_t\Xi_{\bar\theta_{k(t)}}
+O(\sqrt{T\beta_T}).
\]
The $O(\sqrt{T\beta_T})$ absorbs the filter mismatch via~\eqref{eq:dmodel_bound}.

\emph{Optimization.} Exploration cost $=O(\eta\sqrt{T}R_{\max}/(1-\gamma))$. Balancing with term~(A) gives $\eta^\star\sim K_{\mathrm{cov}}^{1/2}\,T^{1/4}(\beta_T/\log T)^{1/4}$ and total $\widetilde O(K_{\mathrm{cov}}^{1/2}T^{3/4}/(1-\gamma)^{3/2})$.
Term~(C) achieves $\widetilde O(\sqrt{T})$; the $T^{3/4}$ rate is driven by the coverage transfer for term~(A).
\end{proof}

% \begin{rmk}[Recovering $\widetilde O(\sqrt{T})$]
% The $T^{3/4}$ rate arises from the $1/\varepsilon$ blow-up in the coverage transfer (Lemma~\ref{lem:coverage}) needed for term~(A). This can be avoided by using an \emph{optimistic} model selection rule: choosing $\tilde\theta_t:=\arg\max_{\theta\in\gC_t}V_\theta^{\star,B}(b_t)$ ensures $\text{(A)}_t\le 0$, leaving only term~(C)$=\widetilde O(\sqrt{T})$. This approach, used in~\cite{liu2023optimistic}, requires solving an optimistic planning problem at each model update, which is computationally more expensive than the $\varepsilon$-greedy approach.
% \end{rmk}
 

%% ===================================================================
\subsection{Optimal Exploration Schedule}
\label{app:exploration_schedule}

Set $\varepsilon_k=\eta/\sqrt{|E_k|}$ so that $\varepsilon_k^2|E_k|=\eta^2$ and $\sum_{j<k}\varepsilon_j^2|E_j|=(k-1)\eta^2$. The exploration-dependent terms scale as:
\begin{align*}
\text{Exploration cost:}&\quad \textstyle\sum_k|E_k|\varepsilon_k=\eta\sum_k\sqrt{|E_k|}=O(\eta\sqrt{T}),\\
\text{Term (A):}&\quad O\bigl(K_{\mathrm{cov}}\sqrt{\beta_T}\cdot T/(\eta\sqrt{\log T}(1-\gamma))\bigr),\\
\text{Term (C):}&\quad O\bigl(\sqrt{T\beta_T}/(1-\gamma)\bigr), \quad\text{independent of }\eta.
\end{align*}
Balancing (A) and exploration:
\[
\eta^\star\sim\left(\frac{K_{\mathrm{cov}}\sqrt{T\beta_T}}{\sqrt{\log T}}\right)^{1/2},
\qquad
\text{(A) + exploration} = \widetilde O\!\left(\frac{K_{\mathrm{cov}}^{1/2}T^{3/4}}{(1-\gamma)}\right).
\]
The total regret is:
\begin{equation}\label{eq:eps_schedule}
\mathrm{Reg}_T = \widetilde O\!\left(\frac{K_{\mathrm{cov}}^{1/2}\,R_{\max}\,T^{3/4}}{(1-\gamma)^{3/2}}\right) + \frac{C_{\mathrm{pf}}\,R_{\max}\,T}{(1-\gamma)^2}\cdot J^{-\frac{1}{2\kappa_{\mathrm{pf}}}} + \frac{C_{\mathrm{mc}}\,R_{\max}\,T}{(1-\gamma)^2\sqrt{N}} + \frac{2\gamma^{d+1}}{(1-\gamma)^2}\sum_t\Xi_{\bar\theta_{k(t)}}.
\end{equation}
The $T^{3/4}$ rate is driven by the $1/\varepsilon$ coverage transfer for term~(A). Using optimistic model selection ($\tilde\theta_t:=\arg\max_{\theta\in\gC_t}V_\theta^{\star,B}$) eliminates term~(A) and recovers $\widetilde O(\sqrt{T})$ from term~(C) alone, at the cost of solving an optimistic planning problem~\citep{liu2023optimistic}.

A constant $\varepsilon$ gives the same $\widetilde O(T^{3/4})$ rate after optimization and may be preferable when $T$ is unknown.



%% ===================================================================

\section{Offline Coverage, Structural Bias Bounds, Case Studies, and Ablations}
\label{app:offline_coverage}
\label{def:spc}
\label{eq:coff_def}
\label{prop:offline_regret}
\label{eq:offline_regret}
\label{eq:hist_l1_hybrid}
\label{eq:A_single_policy}
\label{app:variance_proofs}
\label{aspt:value_embedding}
\label{lem:mean_belief_value}
\label{eq:mean_belief_approx}
\label{eq:upper_step1}
\label{eq:taylor_remainder}
\label{eq:lemma_upper}
\label{eq:lemma_lower}
\label{prop:variance_qmdp_bias}
\label{eq:variance_xi_bound}
\label{eq:tV_upper}
\label{app:case_study_proofs}
\label{cor:lqg}
\label{sec:app08-validation-details}
\label{fig:ch08_bound_scaling}
\label{sec:app08-stress-tests}
\label{fig:ch08_stress_tests}
\label{tab:ch08_active_sensing}
\label{tab:ch08_box_search}
\label{sec:app08-planner-comparison}
\label{tab:ch08_planner_comp}
\label{sec:app08-ant-full}
\label{fig:ch08_ant_full}

\aigen{The greedy-exploitation regret analysis under offline single-policy concentrability (Proposition~\ref{prop:offline_regret}), the posterior-variance curvature bound on the QMDP leaf value bias (Proposition~\ref{prop:variance_qmdp_bias}), the closed-form LQG and contextual MDP case studies (Corollary~\ref{cor:lqg}), and the empirical regret-decomposition audits, stress tests, and depth-$d$ planner ablations appear in Appendices~D.7--D.9 and~E of \citet{tan2026pomdp}.}
